% =====================================================================
%  APPUNTI DI FISICA -- CINEMATICA
%  Compilazione: pdflatex cinematica.tex (due volte, per l'indice)
%  Servono nella stessa cartella: appunti-base.sty e palette-fisica.tex
%  (i disegni sono in TikZ).
% =====================================================================
\documentclass[11pt,a4paper,openany]{report}
\usepackage{../appunti-base}     % stile comune a tutte le materie
% =====================================================================
%  palette-fisica.tex  --  tema della materia FISICA
%  (l'unico file che cambia da una materia all'altra)
% =====================================================================
\newcommand{\MateriaNome}{Fisica 1}
\definecolor{materia}{HTML}{3F5FA8}       % blu (dalla copertina GoodNotes)
\definecolor{materiascuro}{HTML}{2B3A55}  % blu-ardesia
\definecolor{materiachiaro}{HTML}{8FA6D8} % azzurro
\definecolor{accento}{HTML}{C8452F}       % rosso mattone: leggi e principi
\definecolor{esempio}{HTML}{2E8B6E}       % verde: esempi
\definecolor{nota}{HTML}{E0A100}          % ambra: NB
\definecolor{evid}{HTML}{FBE7A1}          % giallo evidenziatore
        % tema della materia (solo colori)

% ---------- Dati della copertina ----------
\universita{Università di Pisa}
\corso{Primo Anno LT in Fisica -- A.A. 2026/2027}
\autore{Cristiano Battini}
\professore{Prof.\ Fabrizio Cei}
\aggiornato{\today}

\begin{document}

\copertina{Cinematica}
\tableofcontents

% ===== da p01.tex =====
% =====================================================================
%  CAPITOLO 1 -- COSA È LA FISICA   (appunti pp. 2-6)
% =====================================================================
\chapter{Cosa è la fisica?}
\lezione{15/09/26}

\section{La fisica e il metodo sperimentale}

Fisica \implica Studia i fenomeni naturali e indotti dall'uomo
\implica Dare una \evid{descrizione quantitativa} basata su principi fondamentali.

\begin{schema}
  \node[nodo] (dati) {Dati sperimentali};
  \node[nodoevid, right=of dati] (teo) {Teorie};
  \node[nodo, above right=0.1cm and 1.1cm of teo] (cost) {Costituenti fondamentali materia};
  \node[nodo, below right=0.1cm and 1.1cm of teo] (int) {Interazioni (forze)};
  \node[nodoevid, below=0.9cm of teo] (pr) {Principi fondamentali};
  \draw[freccia] (dati) -- (teo);
  \draw[freccia] (teo.east) -- (cost.west);
  \draw[freccia] (teo.east) -- (int.west);
  \draw[freccia] (teo) -- (pr);
\end{schema}

\subsection{Metodo sperimentale}
\evid{Metodo sperimentale} \implica Galileo

\begin{schema}
  \node[nodo] (es) {Esperimenti};
  \node[nodo, right=0.8cm of es] (ds) {Dati sperimentali};
  \node[nodo, right=0.8cm of ds] (te) {Teorie};
  \node[nodoevid, below=0.7cm of ds] (tp) {Teoria predittiva};
  \node[nodo, below=0.7cm of tp] (ne) {Nuovi esperimenti};
  \node[nodo, below left=0.8cm and 0.2cm of ne] (conf) {Confermano\\la teoria};
  \node[nodo, below right=0.8cm and 0.2cm of ne] (nconf) {Non confermano\\la teoria};
  \node[nodo, below right=0.7cm and 0.1cm of nconf] (corr) {Correzioni};
  \node[nodo, below left=0.7cm and 0.1cm of nconf] (rif) {Rifiuto};
  \draw[freccia] (es) -- (ds);
  \draw[freccia] (ds) -- (te);
  \draw[freccia] (te) |- (tp);
  \draw[freccia] (tp) -- (ne);
  \draw[freccia] (ne) -- (conf);
  \draw[freccia] (ne) -- (nconf);
  \draw[freccia] (nconf) -- (corr);
  \draw[freccia] (nconf) -- (rif);
  \draw[freccia] (corr.east) to[out=20, in=0] (ne.east);
\end{schema}

\begin{nota}
Un singolo esperimento non basta a confermare una teoria.\\
Ma\dots\ un \chiave{singolo esp.} può smentire una teoria
(se fatto correttamente, si intende).
\end{nota}

\subsection{Quando una teoria va in crisi}
Quando una teoria va in crisi \implica (es.) misure più precise.

Mecc.\ classica \implica Eventi macroscopici in un certo intervallo di $v$ e dim.

Ma\dots
\begin{itemize}
  \item $v \ll c$ \implica Relatività ristretta
  \item $\Delta x \gg$ dim.\ atomiche \implica Mecc.\ quant.
  \item Campi grav.\ molto intensi \implica Relatività gener.
\end{itemize}

Mecc.\ classica quando?
\begin{itemize}
  \item $\Delta x \gg$ atomo
  \item relativistica: $v \ll c$
\end{itemize}

Matematica \implica Strumento.\\
La fisica non è la matematica, ha i dati sperimentali.

% ---------------------------------------------------------------------
\section{Schematizzazione}
\evid{Schematizzazione}
\begin{itemize}
  \item Individuare componenti di un fenomeno \implica effetti dominanti (più significativi)
  \item Trascurare effetti secondari
  \item Includere correzioni
\end{itemize}

\begin{esempio}{caduta gravi}
\begin{itemize}[nosep]
  \item Attrazione grav.\ \implica fondamentale
  \item Resistenza aria
  \item Effetti atmosferici
\end{itemize}
Resistenza aria ed effetti atmosferici: da trascurare. Come? \implica Tubo a vuoto.

Ricreare le condizioni ideali il più possibile \implica Eliminare perturbazioni.

% disegno: p03-tubo (appunti p. 3)
\begin{disegno}
\begin{tikzpicture}[disegno]
  % tubo
  \draw[thick] (0,0) -- (0,2.6);
  \draw[thick] (1,0) -- (1,2.6);
  \draw[thick] (0.5,2.6) ellipse (0.5 and 0.2);
  \draw[thick] (0.5,0.15) ellipse (0.42 and 0.3);
  % piatto e massa
  \draw[thick] (0,2.0) -- (1,2.0);
  \fill (0.5,2.17) ellipse (0.22 and 0.15);
  \draw[->] (0.62,2.25) -- (1.6,2.85) node[right,inner sep=1pt] {$m$};
  \draw (1,2.0) -- (1.6,2.0) node[right] {piatto};
  % pompa
  \draw[thick] (-1.6,0.55) rectangle (-1.0,1.35);
  \fill (-1.4,0.95) circle (1.6pt);
  \fill (0.3,0.95) circle (1.6pt);
  \draw (-1.4,0.95) -- (0.3,0.95);
  \node[anchor=north] at (-1.25,0.5) {pompa};
  % altezza
  \draw[decorate,decoration={brace,amplitude=5pt}] (3.6,2.0) -- (3.6,0.0) node[midway,right=6pt] {$h$};
\end{tikzpicture}
\end{disegno}
\begin{itemize}[nosep]
  \item Misuro $h$ e $t$
  \item Mollo il piatto e osservo la relazione tra $h$ e $t$ \implica $h = \frac12 g t^2$?
\end{itemize}
\end{esempio}

\begin{itemize}
  \item \evid{Fisica fondamentale}: leggi fondamentali della natura \implica principi base
  \item \evid{Fisica applicata}: applicazione della fondamentale
\end{itemize}

\begin{schema}
  \node[nodo] (sc) {Scienza};
  \node[nodo, right=2.2cm of sc] (tec) {Tecnologia};
  \draw[freccia] (sc) to[bend left=25] node[above] {Principi} (tec);
  \draw[freccia] (tec) to[bend left=25] node[below] {Strumenti avanzati \implica es.\ acquisizione dati} (sc);
\end{schema}

\begin{schema}
  \node[nodoevid] (fc) {Fisica Classica};
  \node[below=0.05cm of fc, font=\footnotesize] {Non richiede la mecc.\ quant.\\(e relatività)};
  \node[nodo, above right=1.2cm and 1.4cm of fc, anchor=west] (mec) {Meccanica};
  \node[nodo, right=0.8cm of mec, anchor=west, yshift=0.35cm] (cin) {Cinematica: punto materiale};
  \node[nodo, right=0.8cm of mec, anchor=west, yshift=-0.35cm] (din) {Dinamica: corpi reali (moti + forze)};
  \node[nodo, above right=0.3cm and 1.4cm of fc, anchor=west] (ter) {Termodinamica: calore e lavoro meccanico};
  \node[nodo, below right=0.1cm and 1.4cm of fc, anchor=west] (acu) {Acustica};
  \node[nodo, below right=0.7cm and 1.4cm of fc, anchor=west] (ott) {Ottica};
  \node[nodo, below right=1.3cm and 1.4cm of fc, anchor=west] (ele) {Elettromagnetismo (classico senza relatività)};
  \foreach \n in {mec,ter,acu,ott,ele} \draw[freccia] (fc.east) -- (\n.west);
  \draw[freccia] (mec.east) -- (cin.west);
  \draw[freccia] (mec.east) -- (din.west);
\end{schema}

\begin{schema}
  \node[nodoevid] (fm) {Fisica Moderna};
  \node[nodo, above right=0.1cm and 1.2cm of fm] (rel) {Relatività};
  \node[nodo, below right=0.1cm and 1.2cm of fm] (at) {Atomica e Nucleare};
  \node[nodo, below=0.3cm of at] (mq) {Mecc.\ Quantistica};
  \draw[freccia] (fm.east) -- (rel.west);
  \draw[freccia] (fm.east) -- (at.west);
  \draw[freccia] (at) -- (mq);
\end{schema}

% ---------------------------------------------------------------------
\section{Grandezze fisiche}
\evid{Grandezze fisiche}

Grandezza misurabile:
\begin{itemize}[nosep]
  \item Unità di misura
  \item Tecnica operativa di misura
\end{itemize}

Stessa un.\ di mis.\ \implica \chiave{Omogenee}.

\begin{itemize}
  \item \evid{Misura diretta} \implica Lo strumento misura l'un.\ di mis.
  \item \evid{indiretta} \implica Operazioni matematiche
\end{itemize}

\begin{itemize}
  \item \evid{Grandezze fondamentali} \implica propria unità di misura
  \item \evid{derivate} \implica combinazioni delle fondamentali
\end{itemize}

\subsection{SI -- Sistema unità di misura}
\evid{SI} Sistema unità di misura \implica Sistema internazionale

\begin{center}
\begin{tabular}{lll}
\toprule
\textbf{Grand.} & \textbf{Un.\ mis.} & \textbf{Strum.} \\
\midrule
Lunghezza & metro (m) & Metro \\
Tempo & secondi (s) & Orologio \\
Massa & chilogrammo (kg) & Bilancia e Spettrometro massa \\
Temperatura & Kelvin (K) & Termometro \\
Corrente & Ampere (A) & Amperometro \\
Intens.\ Lum. & candela (cd) & Fotometro \\
\bottomrule
\end{tabular}
\end{center}

Grandezze fondamentali \implica Costanti fisiche da fenomeni naturali.

\begin{esempio}{metro}
$l$ da un rag.\ di luce nel vuoto in $\dfrac{1}{299792458}\,\mathrm{s}$
\end{esempio}
\begin{esempio}{secondo}
Tempo impiegato da $9192631770$ transizioni di una riga atomica del cesio.
\end{esempio}

Rivoluzione della metrologia 2019 \implica Misurazioni più accurate.

Tutte le grandezze fisiche sono ``ancorate'' a costanti fondamentali
\implica Dalla Natura \implica carica elem.; $c$; trans.\ iperfine del cesio, ecc.

\subsection{Multipli e sottomultipli}
\evid{Multipli e Sottomultipli}: si va di mille in mille \implica formalmente.

\begin{center}
\begin{tabular}{rll@{\hspace{2.5em}}rll}
\toprule
$10^{21}$ & zetta & Z & $10^{-3}$ & milli & m \\
$10^{18}$ & exa & E & $10^{-6}$ & micro & $\mu$ \\
$10^{15}$ & peta & P & $10^{-9}$ & nano & n \\
$10^{12}$ & tera & T & $10^{-12}$ & pico & p \\
$10^{9}$ & giga & G & $10^{-15}$ & femto & f \\
$10^{6}$ & mega & M & $10^{-18}$ & atto & a \\
$10^{3}$ & kilo & k & $10^{-21}$ & zepto & z \\
$1$ & & & & & \\
\bottomrule
\end{tabular}
\end{center}

% ---------------------------------------------------------------------
\section{Analisi dimensionale}
\evid{Analisi dimensionale}

\begin{formula}
x = y \iff [x] = [y]
\end{formula}

\lezione{18/9/26}
\subsection{Applicazione analisi dimensionale nel pendolo semplice}

% disegno: p06-pendolo (appunti p. 6)
\begin{disegno}
\begin{tikzpicture}[disegno]
  % soffitto
  \draw[thick] (-0.9,0) -- (1.1,0);
  \foreach \x in {-0.6,-0.2,0.2,0.6} \draw (\x,0.05) -- (\x+0.3,0.35);
  % verticale tratteggiata
  \draw[tratt] (0,0) -- (0,-1.75);
  % filo e massa
  \draw[thick] (0,0) -- (0.55,-1.55) coordinate (B);
  \fill (B) circle (2.2pt);
  \node[right] at (0.22,-0.45) {$\ell$};
  \node at (-0.05,-1.0) [right=0pt] {$\theta$};
  \node[right=2pt] at (B) {$m$};
  % gravita
  \draw[->,thick] (2.0,-0.35) -- (2.0,-1.0) node[midway,right=1pt] {$\vec g$};
\end{tikzpicture}
\end{disegno}

\[
  \tau = \underbrace{l^{\alpha} m^{\beta} g^{\gamma}}_{\text{fattori dimensionali}}
  \ \overbrace{f(\theta)}^{\text{fattore adimensionale}}
\]
\[
  [\tau] = [L]^{\alpha}[m]^{\beta}[g]^{\gamma}, \qquad [g] = [L][T]^{-2}
\]
\[
  [L]^{\alpha+\gamma}[m]^{\beta}[T]^{-2\gamma} = [T] \quad (= [\tau])
\]
\[
  \Rightarrow
  \begin{cases}
    \alpha + \gamma = 0 \\
    \beta = 0 \\
    \gamma = -\frac12
  \end{cases}
  \Rightarrow \alpha = -\gamma \Rightarrow \alpha = \frac12
\]
\begin{formula}
\tau = l \cdot g \cdot f(\theta) = f(\theta)\sqrt{\frac{l}{g}}
\end{formula}
\begin{significato}
Il periodo del pendolo dipende solo dalla lunghezza del filo e da $g$, non dalla massa
($\beta = 0$): un filo più lungo dà oscillazioni più lente, una gravità più forte le rende più rapide.
\end{significato}

\begin{delucidazione}
$f(\theta) = 2\pi$ per piccole osc.\\
Ma questo l'an.\ dim.\ non ce lo dice.
\end{delucidazione}

% ===== da p02.tex =====
% =====================================================================
%  CAPITOLO 2 -- CINEMATICA   (appunti pp. 7-10)
% =====================================================================
\chapter{Cinematica}

Descrivere il moto dei corpi:
\begin{itemize}
  \item Indipendentemente dalle cause che lo producono o modificano
  \item Studio geometrico
\end{itemize}

\begin{delucidazione}
(sul primo punto) quella è la Dinamica
\end{delucidazione}

% ---------------------------------------------------------------------
\section{Punto materiale}
\begin{definizione}{I -- Punto Materiale}
\begin{enumerate}
  \item Oggetto \evid{realmente puntiforme}
        \implica dunque \implica \evid{Dimensioni $\ll$ Spostamento}
  \item Oggetto con \evid{dim $\sim$ spostamento}. Ma:
        \begin{itemize}[nosep]
          \item \evid{non deforma}
          \item \evid{non ruota}
        \end{itemize}
\end{enumerate}
\end{definizione}

\begin{esempio}{}
Terra intorno al sole; treno che va avanti e indietro; auto in frenata.\\
Mentre \implica auto che si ribalta non va bene,
perché le dim sono importanti per la rotazione (ribaltamento).
\end{esempio}

% ---------------------------------------------------------------------
\section{Sistema di riferimento}
\begin{definizione}{II -- Sistema di riferimento}
\evid{Insieme dei corpi rispetto a cui si studia il moto}
\begin{itemize}[nosep]
  \item[\implica] Terna di assi (cartesiani e non)
  \item[\implica] Orologio
  \item[\implica] Unità di misura della lunghezza (metro)
\end{itemize}
\end{definizione}

% disegno: p07-assi (appunti p. 7)
\begin{disegno}
\begin{tikzpicture}[disegno]
  % assi
  \draw[asse] (0,0) -- (0,2.4) node[right] {$z$};
  \draw[asse] (0,0) -- (1.9,0) node[below right] {$y$};
  \draw[asse] (0,0) -- (-0.95,-1.2) node[below right] {$x$};
  % vettore P -> Q (freccia a meta)
  \fill (-0.75,0.25) coordinate (P) circle (1.3pt);
  \fill (1.25,1.05) coordinate (Q) circle (1.3pt);
  \node[above left=-1pt] at (P) {$P$};
  \node[above right=1pt] at (Q) {$Q$};
  \draw[decoration={markings,mark=at position 0.78 with {\arrow{Stealth[length=2.2mm]}}},postaction={decorate}] (P) -- (Q);
  % scritte
  \node[anchor=base] at (4.4,1.3) {$P(x,y,z)$};
  \node[anchor=base] at (4.3,0.45) {$t$};
  \node[anchor=base] at (6.9,1.3) {$Q(x',y',z')$};
  \node[anchor=base] at (7.0,0.45) {$t'$};
  \node[anchor=base] at (5.8,-0.75) {lunghezza fisica};
\end{tikzpicture}
\end{disegno}

% ---------------------------------------------------------------------
\section{Legge oraria e traiettoria}
\begin{definizione}{III -- Legge oraria}
È una \evid{funzione matematica}
\implica Esprime la \evid{dipendenza delle coordinate in funzione del tempo}.
\[
  \underbrace{x(t) \quad y(t) \quad z(t)}_{\text{legge oraria per ogni asse}}
  \ \longrightarrow\ \text{componenti del moto}
\]
\end{definizione}

\begin{definizione}{IV -- Traiettoria}
\evid{Insieme delle posizioni assunte dal corpo durante il suo moto}
\implica Relazione tra coordinate.
\end{definizione}

\begin{nota}
\begin{itemize}[nosep]
  \item Il tempo non compare.
  \item Non è necessariamente una funzione \implica es.\ circ.\ (non è una $f(x)$, Analisi).
\end{itemize}
\end{nota}

% ---------------------------------------------------------------------
\section{Spostamento}
\begin{definizione}{V -- Spostamento}
\[
  t_0 \quad P(x_0, y_0, z_0) \ \longrightarrow\ t_1 \quad Q(x_1, y_1, z_1)
\]
\evid{Differenza (con il suo segno) tra la posizione finale e quella iniziale}
\[
  \text{Spost.} = (\Delta x, \Delta y, \Delta z)
\]
\end{definizione}

\evid{In vettori:}
\begin{align*}
  \vec{r}_0 &: \text{vettore pos.\ a } t_0 \quad (x_0, y_0, z_0)\\
  \vec{r}_1 &: \text{vettore pos.\ a } t_1 \quad (x_1, y_1, z_1)
\end{align*}
\begin{formula}
\Delta\vec{r}\,(t_0 \to t_1) = \vec{r}_1(t_1) - \vec{r}_0(t_0) = (\Delta x, \Delta y, \Delta z)
\end{formula}
\begin{significato}
Lo spostamento è la ``freccia'' che va dal punto di partenza a quello di arrivo:
conta solo dove inizi e dove finisci, non il percorso fatto in mezzo; il segno di ogni
componente dice in che verso ti sei mosso lungo quell'asse.
\end{significato}
\[
  \underset{\vec{r}_i}{\begin{pmatrix} x_i \\ y_i \\ z_i \end{pmatrix}}
  - \underset{\vec{r}_0}{\begin{pmatrix} x_0 \\ y_0 \\ z_0 \end{pmatrix}}
  = \underset{\vec{r}_i - \vec{r}_0}{\begin{pmatrix} x_i - x_0 \\ y_i - y_0 \\ z_i - z_0 \end{pmatrix}}
  = \underset{\Delta\vec{r}}{\begin{pmatrix} \Delta x \\ \Delta y \\ \Delta z \end{pmatrix}}
\]

\begin{esempio}{}
% disegno: p08-esempi (appunti p. 8)
\begin{disegno}
\begin{tikzpicture}[disegno,materia]
  % --- grafico P -> Q ---
  \begin{scope}
    \draw[->,thick] (0,0) -- (3.0,0) node[right] {$x$};
    \draw[->,thick] (0,0) -- (0,1.85) node[above right=-2pt] {$y$};
    \coordinate (P) at (0.85,1.3);
    \coordinate (Q) at (2.6,0.65);
    \draw[dashed] (0,1.3) -- (P) -- (0.85,0);
    \draw[dashed] (0,0.65) -- (Q) -- (2.6,0);
    \draw[thick,fill=white] (P) circle (1.6pt);
    \draw[thick,fill=white] (Q) circle (1.6pt);
    \node[above right=0pt] at (P) {$P$};
    \node[right=2pt] at (Q) {$Q$};
    \draw[decoration={markings,mark=at position 0.6 with {\arrow{Stealth[length=2.2mm]}}},postaction={decorate}]
      ($(P)+(0.25,0.08)$) to[bend left=12] ($(Q)+(-0.1,0.18)$);
    \node[left] at (0,1.3) {$y_0$};
    \node[left] at (0,0.65) {$y_1$};
    \node[below] at (0.85,0) {$x_0$};
    \node[below] at (2.6,0) {$x_1$};
  \end{scope}
  % --- scritte ---
  \node[anchor=base west] at (4.0,1.45) {$P\to Q$};
  \node[anchor=base west] at (4.0,0.85) {$\Delta y<0$};
  \node[anchor=base west] at (4.0,0.25) {$\Delta x>0$};
  \node[anchor=base west] at (6.4,1.45) {es.};
  % --- moto circolare ---
  \begin{scope}[shift={(8.0,-0.4)}]
    \draw[->,thick] (0,0) -- (2.6,-0.05) node[right] {$x$};
    \draw[->,thick] (0,0) -- (0,2.4) node[above right=-2pt] {$y$};
    \draw[thick] (1.6,1.0) ellipse (0.8 and 0.72);
    \draw[thick,-{Stealth[length=2.4mm]}] (1.95,1.67) -- (1.62,1.72);
    \fill[accento] (1.45,1.71) circle (2.2pt);
  \end{scope}
\end{tikzpicture}
\end{disegno}
$P \to Q$: \ $\Delta y < 0$, \ $\Delta x > 0$.\\
Es.\ (moto su una circonferenza): distanza $= 2\pi r$, \ spostamento $= 0$.
\end{esempio}

% ---------------------------------------------------------------------
\newpage
\section{Ascissa curvilinea}
\begin{definizione}{VI -- Ascissa curvilinea $s(t)$}
Equivalente di una coord.\ cartes., ma considerata lungo la traiettoria.
\end{definizione}

% disegno: p09-ascissa (appunti p. 9)
\begin{disegno}
\begin{tikzpicture}[disegno]
  \draw[asse] (-0.5,0) -- (5.0,0) node[right] {$x$};
  \draw[asse] (0,-0.4) -- (0,2.5) node[right] {$y$};
  \node[below left] at (0,0) {$O$};
  \coordinate (O) at (0,0);
  \coordinate (P) at (1.6,1.15);
  \coordinate (Q) at (4.1,2.0);
  \draw[thick,decoration={markings,
      mark=at position 0.15 with {\arrow{Stealth[length=2mm]}},
      mark=at position 0.62 with {\arrow{Stealth[length=2mm]}},
      mark=at position 0.9 with {\arrow{Stealth[length=2mm]}}},postaction={decorate}]
    (O) .. controls (0.3,0.8) and (1.1,0.7) .. (P)
        .. controls (2.1,1.6) and (2.8,1.0) .. (3.4,1.25)
        .. controls (3.8,1.4) and (3.9,1.7) .. (Q);
  \fill (P) circle (1.5pt) node[above=2pt] {$P$};
  \fill (Q) circle (1.5pt) node[above right=0pt] {$Q$};
  % graffe blu
  \draw[materia,thick,decorate,decoration={brace,mirror,amplitude=5pt}]
    (0.05,-0.08) -- (1.55,0.85) node[midway,below right=3pt and 0pt] {$s_0$};
  \draw[materia,thick,decorate,decoration={brace,mirror,amplitude=5pt}]
    (1.85,0.95) -- (4.15,1.55) node[midway,below=6pt] {$s_1$};
\end{tikzpicture}
\end{disegno}

$s_0: O \to P$, \quad $s_1: P \to Q$.

Quindi\dots\ spostamento percorrendo la traiettoria (fisso l'origine $O$): \ $s_0 < 0$, \ $s_1 > 0$.\\
Es.: semicirconferenza da $A$ a $B$ (diametro $2R$): ascissa curvilinea $\pi R$.

\begin{osservazione}
\evid{Spostamento $\neq$ Distanza percorsa}.\\
Andata e ritorno di $r$: \ distanza $= 2r$, \ spostamento $= 0$.
\end{osservazione}

\begin{delucidazione}
Qual è la differenza tra $\vec{x}(t) \to$ spostamento e $\vec{s}(t) \to$ ascissa curvilinea?

% disegno: p09-margine (appunti p. 9)
\begin{disegno}
\begin{tikzpicture}[disegno,
  declare function={f(\x)=1.15+0.32*sin(2.2*(\x-0.9) r)+1.4*(\x>3)*(\x-3)^2;}]
  % ===== Spostamento =====
  \node[accento,anchor=west,font=\normalsize] at (-0.4,2.6) {Spostamento};
  \draw[materia,->,thick] (0,0) -- (4.1,0) node[right] {$t$};
  \draw[materia,->,thick] (0,0) -- (0,2.1) node[right] {$x(t)$};
  \draw[materia,thick,domain=0.65:3.55,smooth,samples=60] plot (\x,{f(\x)});
  \coordinate (A) at (1.0,{f(1.0)});
  \coordinate (B) at (2.95,{f(2.95)});
  \draw[accento,thick] (A) -- (B);
  \fill[accento] (A) circle (2pt) (B) circle (2pt);
  \draw[accento,dashed] (A) -- (1.0,0) node[below] {$t_0$};
  \draw[accento,dashed] (B) -- (2.95,0) node[below] {$t$};
  \node[accento] at (2.4,1.85) {$x(t_0,t)$};
  % ===== Ascissa curvilinea =====
  \begin{scope}[yshift=-3.5cm]
  \node[esempio,anchor=west,font=\normalsize] at (-0.8,2.6) {Ascissa curvilinea};
  \draw[materia,->,thick] (0,0) -- (4.1,0) node[right] {$t$};
  \draw[materia,->,thick] (0,0) -- (0,2.1) node[right] {$s(t)$};
  \draw[materia,thick,domain=0.65:3.55,smooth,samples=60] plot (\x,{f(\x)});
  \foreach \a/\b in {0.75/1.05,1.05/1.55,1.55/2.0,2.0/2.45,2.45/2.85,2.85/3.3}
    \draw[esempio,very thick,-{Stealth[length=1.8mm]}] (\a,{f(\a)}) -- (\b,{f(\b)});
  \draw[esempio,dashed] (0.75,{f(0.75)}) -- (0.75,0) node[below] {$t_0$};
  \draw[esempio,dashed] (3.3,{f(3.3)}) -- (3.3,0) node[below] {$t$};
  \node[esempio] at (0.55,1.55) {$s_0$};
  \node[esempio] at (1.2,1.75) {$s_1$};
  \node[esempio] at (1.95,1.65) {$s_2$};
  \node[esempio] at (2.4,1.55) {$s_3$};
  \node[esempio] at (2.8,1.55) {$s_4$};
  \node[esempio] at (3.2,1.65) {$s_5$};
  \node[esempio] at (2.7,2.15) {$s(t,t_0)$};
  \end{scope}
\end{tikzpicture}
\end{disegno}

\textbf{Spostamento}: lunghezza descritta dalla retta che collega i due punti.\\
\textbf{Ascissa curvilinea}: lunghezza descritta dalla curva che collega due punti.
\end{delucidazione}

% ---------------------------------------------------------------------
\section{Relazione legge oraria \texorpdfstring{$\to$}{->} traiettoria}
\evid{Relazione legge oraria $\to$ traiettoria} (l'inverso non è possibile).

\evid{Tecnica standard}: ricavare il tempo da una legge oraria e sostituirla nelle altre.

\begin{esempio}{moto di gravi lanciati}
Legge oraria:
\[
  \begin{cases}
    x(t) = v_x t \\
    y(t) = v_y t - \frac12 g t^2
  \end{cases}
  \qquad
  t = \frac{x}{v_x} \ \longmapsto\
  \underbrace{y(x) = \frac{v_y}{v_x}\,x - \frac{g}{2}\left(\frac{x}{v_x}\right)^2}_{\text{traiettoria}}
\]
\implica parabola (termine quadratico).

$\bar{x}$: $x$ a cui $y$ è $0$:
\[
  y(\bar{x}) = 0 \ \longmapsto\ \bar{x}\,\frac{v_y}{v_x} - \frac{g}{2}\left(\frac{\bar{x}}{v_x}\right)^2 = 0
\]
\[
  \bar{x}\left(\frac{v_y}{v_x} - \frac{g}{2v_x^2}\,\bar{x}\right) = 0
  \qquad \bar{x} = 0
  \qquad \bar{x} = -\frac{v_y}{v_x}\left(-\frac{2v_x^2}{g}\right) = \frac{2v_xv_y}{g}
\]
% disegno: p09-parabola (appunti p. 9)
\begin{disegno}
\begin{tikzpicture}[disegno,materia]
  \draw[->,thick] (0,0) -- (4.3,0) node[right] {$x$};
  \draw[->,thick] (0,0) -- (0,3.0) node[right=2pt] {$y$};
  \coordinate (V) at (1.95,2.05);
  \draw[thick] (0,0) .. controls (0.5,1.6) and (1.3,2.05) .. (V)
                 .. controls (2.6,2.05) and (3.7,1.2) .. (3.75,0);
  \draw[dashed] (0,2.05) -- (V) -- (1.95,0);
  \fill (V) circle (2pt) node[above=2pt] {$V$};
  \fill (0,0) circle (1.3pt);
  \node[below] at (0,-0.05) {$O$};
  \node[below] at (3.75,0) {$\bar x$};
\end{tikzpicture}
\end{disegno}
\end{esempio}

\begin{esempio}{in cui non funziona}
\[
  \begin{cases}
    x = R\cos(\omega t) \ \to\ \dfrac{x}{R} = \cos(\omega t) \ \to\ t = \dfrac{1}{\omega}\arccos\!\left(\dfrac{x}{R}\right)\\[1ex]
    y = R\sin(\omega t)
  \end{cases}
\]
\[
  y(x) = R\sin\!\left(\omega\,\frac{1}{\omega}\arccos\!\left(\frac{x}{R}\right)\right)
\]
Concettualmente errata \implica non rispetta l'arcocoseno.

\begin{delucidazione}
\aggiunto\ L'arcocoseno restituisce solo angoli tra $0$ e $\pi$, quindi
$t = \frac{1}{\omega}\arccos\!\left(\frac{x}{R}\right)$ vale solo per mezzo giro.
Sostituendo si ottiene $y = R\sin\!\left(\arccos\frac{x}{R}\right) = R\sqrt{1 - \frac{x^2}{R^2}} \geq 0$:
solo la metà superiore della circonferenza.
Il moto circolare è periodico: a ogni $x$ corrispondono due punti ($y > 0$ e $y < 0$) e infiniti istanti,
quindi $t$ non si ricava in modo unico da $x$. Per questo il tempo si elimina con
$\cos^2 + \sin^2 = 1$, come subito dopo.
\end{delucidazione}

Invece:
\[
  \begin{cases}
    \cos(\omega t) = \dfrac{x}{R} \\[1ex]
    \sin(\omega t) = \dfrac{y}{R}
  \end{cases}
  \ \to\ \cos^2(\omega t) + \sin^2(\omega t) = \frac{x^2}{R^2} + \frac{y^2}{R^2} = 1
\]
\[
  x^2 + y^2 = R^2 \ \to\ \text{circonf.}
\]
Se i coef.\ di $x$ e $y$ non sono uguali, $a \neq b$:
\[
  \begin{cases}
    x = a\cos(\omega t) \\
    y = b\sin(\omega t)
  \end{cases}
  \quad
  \underbrace{\cos^2(\omega t) + \sin^2(\omega t)}_{1} = \frac{x^2}{a^2} + \frac{y^2}{b^2}
\]
\[
  x^2 + y^2 = a^2 + b^2 \ : \ \text{ellisse}
\]
\end{esempio}

\begin{delucidazione}
\textbf{Esercizio}:
$\begin{cases} x = a\cosh(\omega t) \\ y = b\sinh(\omega t) \end{cases}$
\qquad
$\cosh x = \dfrac{e^x + e^{-x}}{2}$, \quad $\sinh x = \dfrac{e^x - e^{-x}}{2}$
\end{delucidazione}

% ===== da p03.tex =====
% =====================================================================
%  CAPITOLO 3 -- MOTO UNIDIMENSIONALE   (appunti pp. 10-13)
% =====================================================================
\chapter{Moto unidimensionale}

(o rettilineo)

Vuol dire \implica suffic.\ \chiave{1 sola coordinata} (non sempre cartesiana)
\implica 1 sola legge oraria.

\begin{esempio}{}
Il moto circolare è unidimensionale \implica mi serve solo l'angolo $\theta(t)$.
\end{esempio}

\nonewpage
\section{Moto rettilineo}
\evid{Moto rettilineo} \implica Moto unidim., con \evid{traiettoria $=$ segmento}.

% disegno: p10-retta (appunti p. 10)
\begin{disegno}
\begin{tikzpicture}[disegno]
  \draw[asse,black] (0,0) -- (5.6,0) node[right] {$x$};
  \foreach \x/\i in {1.3/0,2.9/1,4.4/2} {
    \draw[fill=white] (\x,0) circle (1.6pt);
    \node[above=2pt] at (\x,0) {$t_\i$};
    \node[below=2pt] at (\x,0) {$x_\i$};
  }
\end{tikzpicture}
\end{disegno}

\subsection{Rappresentazione grafica di \texorpdfstring{$x(t)$}{x(t)}}
\evid{Rappresentazione graf.~$x(t)$ legge oraria}

% disegno: p10-grafico (appunti p. 10)
\begin{disegno}
\begin{tikzpicture}[disegno,x=1.15cm,y=1.15cm]
  \draw[asse] (0,0) -- (4.6,0) node[right] {$t$};
  \draw[asse] (0,0) -- (0,2.0) node[above right=-3pt and -2pt] {$x(t)$};
  \fill (0,0) circle (1.4pt);
  \draw[thick] (0,0) .. controls (0.35,0.5) and (0.5,1.2) .. (0.85,1.2)
     .. controls (1.15,1.2) and (1.3,0.95) .. (1.7,0.95)
     .. controls (2.1,0.95) and (2.25,1.3) .. (2.35,1.6)
     .. controls (2.45,1.9) and (2.7,1.9) .. (2.95,1.85)
     .. controls (3.05,1.82) and (3.1,1.75) .. (3.2,1.8)
     .. controls (3.4,1.85) and (3.4,1.4) .. (3.55,1.1)
     .. controls (3.65,0.9) and (3.75,0.85) .. (3.9,0.75);
  \draw[tratt,black] (0.8,1.2) -- (0.8,0);
  \draw[tratt,black] (1.65,0.95) -- (1.65,0);
  \draw[tratt,black] (2.6,1.88) -- (2.6,0);
  \node[below] at (0,0) {$t_0$};
  \node[below] at (0.8,0) {$t_1$};
  \node[below] at (1.65,0) {$t_2$};
  \node[below] at (2.6,0) {$t_3$};
  % commento grigio
  \draw[black!55,thick,decorate,decoration={brace,mirror,amplitude=5pt}] (-0.4,-0.55) -- (2.95,-0.55);
  \node[black!55,anchor=west] at (-0.5,-1.15) {Come varia $x(t)$};
\end{tikzpicture}
\end{disegno}

\begin{nota}
Non è la traiettoria! Quella è un segmento: il grafico mostra come varia $x(t)$.\\
Può essere negativa!
\end{nota}

\begin{esempio}{1}
\[
  x(t) = x_0 + \overbrace{v(t - t_0)}^{[L]} \qquad v > 0
\]
% disegno: p11-es1 (appunti p. 11)
\begin{disegno}
\begin{tikzpicture}[disegno,materia]
  \draw[->,thick] (-0.35,0) -- (4.0,0) node[below right=-1pt] {$t$};
  \draw[->,thick] (0,-0.5) -- (0,2.5) node[above] {$x$};
  \coordinate (A) at (0.95,0.8);
  \coordinate (B) at (3.6,2.2);
  \draw[dashed] (0,0.8) -- (A) -- (0.95,0);
  \draw[thick] (A) to[bend left=8] (B);
  \fill (B) circle (1.8pt);
  \node[left] at (0,0.85) {$x_0$};
  \node[below] at (0.95,0) {$t_0$};
\end{tikzpicture}
\end{disegno}
Il punto materiale percorre spazi uguali in\dots\ tempi uguali.
\[
  t_1, t_2 \to x(t_1), x(t_2) \qquad t_2 > t_1
\]
\[
  \frac{x(t_1) + x(t_2)}{t_1 - t_2}
  = \frac{\overbrace{v(t_2 - t_1)}^{\text{spazio}}}{\underbrace{t_2 - t_1}_{\text{tempo}}}
  = v \ \to\ \text{cost.}
\]
\end{esempio}

\sinewpage

\begin{esempio}{2}
\[
  y(t) = y_0 - a(t - t_0)^2 \qquad t \geq t_0
\]
% disegno: p11-es2 (appunti p. 11)
\begin{disegno}
\begin{tikzpicture}[disegno]
  \draw[asse,materia] (0,0) -- (4.2,0) node[below right] {$t$};
  \draw[asse,materia] (0,0) -- (0,2.6) node[above right] {$y$};
  \draw[tratt,materia] (0,1.9) node[left,text=materia] {$y_0$} -- (0.85,1.9);
  \draw[tratt,materia] (0.85,1.9) -- (0.85,0) node[below,text=materia] {$t_0$};
  \draw[curva] (0.85,1.9) parabola (2.5,1.05);
  \draw[materia,thick,fill=white] (2.5,1.05) circle (1.6pt);
\end{tikzpicture}
\end{disegno}
\begin{itemize}[nosep]
  \item parabola
  \item la traiettoria è comunque un segmento
\end{itemize}
\end{esempio}

\begin{esempio}{3: incrocio tra leggi orarie}
% disegno: p11-es3 (appunti p. 11)
\begin{disegno}
\begin{tikzpicture}[disegno]
  \draw[asse,materia] (0,0) -- (4.4,0) node[below right] {$t$};
  \draw[asse,materia] (0,0) -- (0,3.6) node[above] {$x$};
  % curva crescente (concava, flesso orizzontale all'intersezione, poi risale)
  \draw[curva] (0.55,0.35) .. controls (0.9,1.3) and (1.7,1.65) .. (2.6,1.65)
               .. controls (3.1,1.65) and (3.35,1.7) .. (3.8,2.3);
  % curva decrescente (inizio in alto, flesso, minimo vicino all'asse)
  \draw[curva] (1.45,2.6) .. controls (2.0,2.5) and (2.3,2.1) .. (2.6,1.65)
               .. controls (2.95,1.0) and (3.2,0.4) .. (3.7,0.35)
               .. controls (3.9,0.33) and (4.05,0.38) .. (4.1,0.45);
  \node[punto,materia] at (2.6,1.65) {};
  \draw[tratt,materia] (0,1.65) node[left,text=materia] {$\bar{x}$} -- (2.6,1.65);
  \draw[tratt,materia] (2.6,1.65) -- (2.6,0) node[below,text=materia] {$\bar{t}$};
\end{tikzpicture}
\end{disegno}
$\bar{x}, \bar{t}$: radici eq.:
\[
  x_1(t) = x_2(t) \ \Longrightarrow\ t = \bar{t}
\]
\end{esempio}

% ---------------------------------------------------------------------
\section{Velocità}
\begin{definizione}{7 -- Velocità media}
\evid{Velocità media}: rapporto tra la differenza delle posizioni assunte
nell'istante $t_1$ e $t_0$ e la differenza $(t_1 - t_0)$ \quad $(t_1 > t_0)$.
\end{definizione}
\begin{formula}
v_m(t_0, t_1) = \frac{x_1 - x_0}{t_1 - t_0} = \frac{x(t_1) - x(t_0)}{t_1 - t_0}
\end{formula}
\begin{significato}
Dice di quanto è cambiata in media la posizione per ogni unità di tempo
nell'intervallo considerato: non dice nulla su cosa è successo in mezzo
(il corpo può aver accelerato, frenato o essere tornato indietro).
\end{significato}

\begin{definizione}{8 -- Velocità istantanea $v(t_0)$}
$\Delta t \to 0$, \quad $t_1 = t_0 + \Delta t$
\[
  v_m = \frac{x(t_0 + \Delta t) - x(t_0)}{\Delta t}
\]
La velocità all'ist.\ $t_0$ è:
\[
  v(t_0) = \lim_{\Delta t \to 0} v_m(t_0, t_0 + \Delta t)
  = \lim_{\Delta t \to 0} \underbrace{\frac{x(t_0 + \Delta t) - x(t_0)}{\Delta t}}_{\text{rapporto incrementale } \to \text{ derivata}}
\]
\end{definizione}
\begin{formula}
v(t_0) = \frac{d}{dt}\,x(t) = \underbrace{\dot{x}(t)}_{\text{notazione derivata}}
\end{formula}
\begin{significato}
È la velocità ``in quell'istante preciso'': si prende la velocità media su intervalli
di tempo sempre più piccoli attorno a $t_0$. Misura quanto rapidamente sta cambiando
la posizione proprio in quel momento.
\end{significato}
\[
  [v_m] = [v] = \frac{[L]}{[T]} = \frac{\mathrm{m}}{\mathrm{s}}
\]

\subsection{Interpretazione geometrica della velocità}
\evid{Interpretazione geometrica della velocità}

\needspace{10\baselineskip}
\evid{Vel.~media}:
% disegno: p12-media (appunti p. 12)
\begin{disegno}
\begin{tikzpicture}[disegno]
  \draw[asse] (0,-0.05) -- (7.0,-0.05) node[below right] {$t$};
  \draw[asse] (0,-0.15) -- (0,2.6) node[above right] {$x(t)$};
  % tratteggi
  \draw[tratt,black] (0,1.75) node[left] {$x_1$} -- (1.8,1.75);
  \draw[tratt,black] (0,1.3) node[left] {$x_0$} -- (3.25,1.3);
  \draw[tratt,black] (1.56,0.75) -- (1.56,-0.05) node[below] {$t_0$};
  \draw[tratt,black] (2.12,1.3) -- (2.12,-0.05) node[below] {$t_1$};
  % curva (nera)
  \draw[thick,black] (0.69,0.31) .. controls (0.7,0.8) and (0.95,1.08) .. (1.31,1.1);
  \draw[thick,black] (1.8,1.75) .. controls (2.6,2.0) and (3.2,2.45) .. (3.56,2.31)
        .. controls (3.85,2.2) and (3.95,1.8) .. (4.25,1.75);
  % secante
  \draw[thick,materia] (0.98,0.55) -- (2.5,2.75) node[right,yshift=2pt] {secante};
  % angolo
  \draw[thick,nota] (1.8,1.75) ++(-45:0.38) arc[start angle=-45,end angle=15,radius=0.38];
  \node[text=nota] at (2.42,1.62) {$\alpha$};
\end{tikzpicture}
\end{disegno}
\[
  v_m = \frac{x_1 - x_0}{t_1 - t_0} = \tan\alpha
\]
$v_m \to$ coef.\ ang.\ secante $(t_0, x_0)\,(t_1, x_1)$ (rapporto incr.).

\needspace{10\baselineskip}
\evid{Vel.~ist.}:
% disegno: p12-ist (appunti p. 12)
\begin{disegno}
\begin{tikzpicture}[disegno]
  \draw[asse] (-0.25,0) -- (5.8,0);
  \draw[asse] (0,-0.2) -- (0,3.0);
  % tratteggi
  \draw[tratt,black] (0,1.3) node[left] {$x_0$} -- (1.43,1.3);
  \draw[tratt,black] (1.57,1.15) -- (1.57,0) node[below,xshift=-2pt] {$t_0$};
  \draw[tratt,accento] (0,1.6) node[left,yshift=3pt] {$x_1$} -- (1.77,1.6) -- (1.77,0) node[below,xshift=3pt] {$t_1$};
  % curva
  \draw[thick,black] (0.64,0.67) .. controls (0.66,1.25) and (1.0,1.58) .. (1.57,1.58)
        .. controls (2.3,1.58) and (3.3,1.55) .. (3.71,2.57);
  % tangente
  \draw[very thick,esempio] (-0.3,2.04) -- (3.36,1.14) node[right=5pt,text=esempio] {tang};
  \draw[thick,esempio,->] (5.75,0.8) .. controls (6.25,0.9) and (6.0,-0.2) .. (6.35,-0.85);
\end{tikzpicture}
\end{disegno}
$t_1 \to t_0$, \quad $x_1 \to x_0$\\
\implica La secante diventa la tang.\\
$v \to$ coef.\ ang.\ tang.\ (derivata).

\subsection{Quando esiste la velocità istantanea}
Vel.\ istantanea? \implica $x(t)$ continua e derivabile.

\begin{esempio}{}
Se un pallone rimbalza sul muro ($t = 0^-$ va verso il muro, $t = 0^+$ torna indietro, urto a $t = 0$):\\
a $t = 0$ \chiave{non posso calcolare $v$} \implica Calcolo $v(0^+)$ e $v(0^-)$.
\end{esempio}

% disegno: p13-segno (appunti p. 13)
\begin{disegno}
\begin{tikzpicture}[disegno,x=0.95cm,y=0.95cm]
  \draw[asse] (0,-1.75) -- (0,2.7) node[above right] {$x(t)$};
  \draw[asse] (0,0) -- (3.3,0) node[right] {$t$};
  \draw[thick,black] (0,0) .. controls (0.4,0.7) and (0.8,0.93) .. (1.07,0.93)
     .. controls (1.5,0.93) and (1.8,0.45) .. (2.0,0)
     .. controls (2.25,-0.5) and (2.5,-0.85) .. (2.7,-0.88);
\end{tikzpicture}
\end{disegno}
\[
  v = \frac{d}{dt}\,x(t)
  \begin{cases}
    > 0 & x \text{ cresce} \\
    < 0 & x \text{ decresce}
  \end{cases}
\]

% ===== da p04.tex =====
% =====================================================================
%  CAPITOLO 4 -- MOTI RETTILINEI: MRU E MRUA   (appunti pp. 13-17)
% =====================================================================
\chapter{Moto rettilineo uniforme e uniformemente accelerato}

\section{Moto rettilineo uniforme (MRU)}
Moto in cui $v = \evid{v_m = \text{costante}} = v_c$.

\begin{formula}
x(t) = x_0 + v_c(t - t_0)
\end{formula}
\begin{significato}
La posizione cambia della stessa quantità in tempi uguali: partendo da $x_0$
all'istante $t_0$, ogni secondo si aggiungono $v_c$ metri. Il grafico $x(t)$ è una retta
la cui pendenza è la velocità (salita se $v_c > 0$, discesa se $v_c < 0$).
\end{significato}
\[
  \Longrightarrow
  \begin{cases}
    v_c = \dfrac{x(t) - x_0}{t - t_0} \\[1.5ex]
    t = t_0 + \dfrac{x(t) - x_0}{v_c}
  \end{cases}
\]

Grafico:
% disegno: p13-mru (appunti p. 13)
\begin{disegno}
\begin{tikzpicture}[disegno,x=0.85cm,y=0.85cm]
  % grafico x(t) (nero)
  \draw[asse] (0,-1.26) -- (0,1.6) node[above right] {$x(t)$};
  \draw[asse] (-0.83,0) -- (3.75,0) node[below right] {$t$};
  \draw[tratt,black] (0,0.5) node[left] {$x_0$} -- (1.06,0.5);
  \draw[black!80] (1.02,0.36) -- (1.01,0.24);
  \node[below] at (1.1,0) {$t_0$};
  \draw[thick,black] (1.06,0.5) -- (3.22,1.19);
  \draw[thick,black] (1.06,0.5) -- (2.11,-0.53);
  % annotazioni (grigio)
  \begin{scope}[black!55]
    \node (p) at (3.95,1.13) {$v_c>0$};
    \draw[thick] (p) ellipse (0.68 and 0.42);
    \node (n) at (2.95,-0.75) {$v_c<0$};
    \draw[thick] (n) ellipse (0.68 and 0.42);
    \draw[->,thick] (4.65,1.48) -- (5.65,1.62);
    \draw[->,thick] (3.45,-1.0) to[bend right=15] (4.15,-1.33);
    % v(t) positivo
    \draw[->,thick] (7.06,0.41) -- (7.06,2.47) node[above right] {$v(t)$};
    \draw[->,thick] (6.67,0.78) -- (9.25,0.78) node[below right] {$t$};
    \draw[thick] (6.67,1.21) -- (9.45,1.21) node[right] {positivo};
    % v(t) negativo
    \draw[->,thick] (4.89,-1.76) -- (4.89,-0.26) node[above right] {$v(t)$};
    \draw[->,thick] (4.44,-0.76) -- (7.1,-0.76) node[right] {$t$};
    \draw[thick] (4.44,-1.33) -- (7.3,-1.33) node[right] {negativo};
  \end{scope}
\end{tikzpicture}
\end{disegno}
$v_c > 0$: $v(t)$ positivo; \quad $v_c < 0$: $v(t)$ negativo.

\begin{nota}
È \evid{fondamentale per i sistemi di riferimento inerziali}.
\end{nota}

% ---------------------------------------------------------------------
\section{Moto rettilineo non uniforme}
\evid{Moto rett.~non uniforme} \implica $v = v(t)$: la velocità dipende dal tempo
\implica grafico $v(t), t$.

% disegno: p13-nonunif (appunti p. 13)
\begin{disegno}
\begin{tikzpicture}[disegno,x=0.8cm,y=0.8cm]
  % ---- grafico x(t) ----
  \draw[asse] (-0.4,0) -- (3.2,0) node[right] {$t$};
  \draw[asse] (0,-0.6) -- (0,2.1) node[right,xshift=2pt] {$x(t)$};
  \draw[tratt,black] (0.53,0) -- (0.53,0.5);
  \draw[tratt,black] (1.44,0) -- (1.44,0.91);
  \draw[tratt,black] (1.96,0) -- (2.15,1.55);
  \draw[tratt,black] (2.34,0) -- (2.36,0.75);
  \draw[thick,black] (0.53,0.5) .. controls (0.9,0.52) and (1.25,0.7) .. (1.44,0.91)
     .. controls (1.75,1.3) and (1.8,1.92) .. (2.12,1.92)
     .. controls (2.5,1.92) and (2.38,1.3) .. (2.42,0.97)
     .. controls (2.44,0.5) and (2.5,0.37) .. (2.72,0.37)
     .. controls (2.92,0.37) and (3.0,0.5) .. (3.08,0.69);
  \draw[very thick,black] (2.56,0.37) -- (2.88,0.37);
  \foreach \p in {(0.53,0.5),(1.44,0.91),(2.2,1.8),(2.42,0.97)}
    \node[punto,inner sep=0.9pt] at \p {};
  \node[below] at (0.53,0) {$t_0$};
  \node[below] at (1.38,0) {$t_1$};
  \node[below] at (1.94,0) {$t_2$};
  \node[below] at (2.5,0) {$t_3$};
  % ---- annotazioni grigie ----
  \begin{scope}[black!55]
    \draw[thick] (2.72,0.36) ellipse (0.42 and 0.2);
    \draw[->,thick] (3.15,0.38) -- (3.72,0.59);
    \node[right] at (3.75,0.72) {vel${}=0$};
    \draw[thick,decorate,decoration={brace,mirror,amplitude=5pt}] (4.05,0.25) -- (5.95,0.25);
    \node[align=center,below] at (5.15,0.0) {punto a\\tangenza\\orizzontale};
    \draw[thick,decorate,decoration={brace,mirror,amplitude=5pt}] (-0.05,-1.6) -- (2.45,-1.6);
    \node at (1.4,-2.35) {$x(t)$};
    \node at (7.06,1.11) {$\longmapsto$};
    \draw[thick,decorate,decoration={brace,mirror,amplitude=5pt}] (9.15,-0.85) -- (12.4,-0.85);
    \node at (10.8,-1.75) {$v(t)$};
  \end{scope}
  % ---- grafico v(t) ----
  \begin{scope}[shift={(8.9,0.5)},xscale=1.25]
    \draw[asse] (-0.3,0) -- (2.6,0) node[right] {$t$};
    \draw[asse] (0,-0.4) -- (0,2.0) node[above right] {$v(t)$};
    \draw[tratt,black] (0.51,0) -- (0.51,0.56);
    \draw[tratt,black] (0.98,0) -- (0.98,0.94);
    \draw[black!80] (1.92,-0.08) -- (1.92,0.1);
    \draw[tratt,black] (1.95,0) -- (2.0,-0.85);
    \draw[very thick,black] (0.49,0.56) .. controls (0.55,0.85) and (0.75,0.94) .. (0.97,0.94)
       .. controls (1.15,0.94) and (1.3,0.6) .. (1.49,0)
       .. controls (1.7,-0.55) and (1.85,-1.05) .. (2.21,-0.95);
    \node[below] at (0.47,0) {$t_0$};
    \node[below] at (0.98,0) {$t_1$};
    \node[below right,xshift=-3pt] at (1.5,0) {$t_2$};
    \node[below right,xshift=-1pt] at (1.95,0) {$t_3$};
  \end{scope}
  % ---- commento ----
  \draw[thick] (13.95,2.4) .. controls (14.2,2.4) and (14.0,1.4) .. (14.22,1.3)
     .. controls (14.0,1.2) and (14.2,0.2) .. (13.95,0.2);
  \node[right,align=left] at (14.25,1.3) {non\\uniforme};
\end{tikzpicture}
\end{disegno}
Nei punti a tangenza orizzontale di $x(t)$: vel.\ $= 0$. \quad Non uniforme.

% ---------------------------------------------------------------------
\section{Moto rettilineo uniformemente accelerato (MRUA)}
\evid{Accelerazione $=$ costante}.

\subsection{Accelerazione media}
\begin{definizione}{Accelerazione media}
\[
  a_m(t_0, t_1) = \frac{v(t_1) - v(t_0)}{t_1 - t_0}
\]
\end{definizione}
\begin{significato}
Dice di quanto è cambiata in media la velocità per ogni unità di tempo:
come la velocità media è la pendenza della secante del grafico $x(t)$,
l'accelerazione media è la pendenza della secante del grafico $v(t)$.
\end{significato}
% disegno: p14-am (appunti p. 14)
\begin{disegno}
\begin{tikzpicture}[disegno,x=0.85cm,y=0.85cm]
  \draw[asse] (-0.6,0) -- (3.6,0) node[right] {$t$};
  \draw[asse] (0,-0.6) -- (0,3.0) node[above right] {$v(t)$};
  \draw[tratt,black] (0,2.14) node[left] {$v_2$} -- (2.96,2.14) -- (2.96,0) node[below] {$t_2$};
  \draw[tratt,black] (0,1.22) node[left] {$v_1$} -- (0.89,1.22) -- (0.89,0) node[below] {$t_1$};
  \draw[thick,black] (0.6,0.55) .. controls (0.65,0.95) and (0.75,1.15) .. (0.89,1.22)
      .. controls (1.2,1.42) and (1.5,1.5) .. (1.75,1.52)
      .. controls (2.3,1.6) and (2.7,1.8) .. (3.0,2.2);
  \draw[thick,nota] (0.14,0.82) -- (3.7,2.43);
  \draw[thick,dashed,nota] (0.96,1.13) -- (3.8,1.17);
  \draw[->,black] (3.85,2.33) -- (4.5,2.33);
  \node[anchor=north west,align=left,inner sep=1pt] at (4.55,2.58) {coefficiente\\della retta\\\quad secante};
\end{tikzpicture}
\end{disegno}
\implica coefficiente della retta secante.
\[
  [a_m] = \frac{[v]}{[T]} = \frac{[L]}{[T]}\cdot\frac{1}{[T]} = \frac{[L]}{[T^2]} \ \to\ \frac{\mathrm{m}}{\mathrm{s}^2}
\]

\subsection{Accelerazione istantanea}
\begin{definizione}{Accelerazione istantanea}
\[
  a(t_0) = \lim_{\Delta t \to 0} a_m
  = \underbrace{\left.\frac{d}{dt}\,v(t)\right|_{t_0} = \dot{v}(t_0)}_{\text{derivata della velocità}}
  = \underbrace{\frac{d^2}{dt^2}\,x = \ddot{x}(t_0)}_{\text{derivata seconda della posizione}}
\]
\end{definizione}
\begin{significato}
Misura quanto rapidamente sta cambiando la velocità in quell'istante:
è la pendenza della tangente al grafico $v(t)$, cioè la ``curvatura'' del grafico $x(t)$.
\end{significato}
$a_t$ è la tangente alla curva in $(t, v(t))$ dell'accelerazione.

\subsection{Accelerazione costante}
Acceler.\ costante ($> 0$) $= a_c$:
\begin{formula}
v(t) = v(t_0) + a(t - t_0)
\end{formula}
\begin{significato}
Con accelerazione costante la velocità cresce (o cala) della stessa quantità in tempi uguali:
il grafico $v(t)$ è una retta di pendenza $a_c$.
\end{significato}
\[
  a_c = \frac{v(t) - v(t_0)}{t - t_0}
  \qquad
  t = t_0 + \left(\frac{v(t) - v(t_0)}{a_c}\right)
\]
\evid{La velocità è una retta.}

% disegno: p14-grafici (appunti p. 14)
\begin{disegno}
\begin{tikzpicture}[disegno,x=0.0087cm,y=-0.0087cm]
  % (coordinate in pixel dell'originale)
  % --- v(t), a>0 ---
  \draw[asse] (155,375) -- (155,130) node[right] {$v(t)$};
  \draw[asse] (118,338) -- (395,338) node[below right] {$t$};
  \draw[tratt,black] (155,268) -- (225,268) -- (225,338);
  \draw[thick,black] (225,268) -- (360,178);
  % --- a(t), a>0 ---
  \draw[asse] (633,375) -- (633,135) node[right] {$a(t)$};
  \draw[asse] (603,338) -- (880,338) node[below right] {$t$};
  \draw[tratt,black] (633,265) -- (695,265) -- (695,338);
  \draw[tratt,black] (830,265) -- (830,338);
  \draw[very thick,black] (695,265) -- (830,265);
  \node[right] at (900,180) {$a>0 \quad a$ costante};
  % --- v(t), a<0 ---
  \draw[asse] (158,742) -- (158,525) node[above,xshift=4pt] {$v(t)$};
  \draw[asse] (138,742) -- (415,742) node[below right] {$t$};
  \draw[tratt,black] (158,601) -- (230,601) -- (230,540);
  \draw[thick,black] (230,601) -- (358,718);
  % --- a(t), a<0 ---
  \draw[asse] (633,742) -- (633,525) node[above right,xshift=-4pt] {$a(t)$};
  \draw[asse] (615,601) -- (845,601) node[right] {$t$};
  \draw[tratt,black] (633,668) -- (695,668) -- (695,601);
  \draw[tratt,black] (828,668) -- (828,601);
  \draw[very thick,black] (695,668) -- (828,668);
  \node[right] at (965,635) {$a<0 \quad a$ costante};
\end{tikzpicture}
\end{disegno}
$a > 0$: $a$ costante; \quad $a < 0$: $a$ costante.

\subsection{Ricavare \texorpdfstring{$x(t)$}{x(t)} da \texorpdfstring{$v(t_0)$}{v(t0)} e \texorpdfstring{$x(t_0)$}{x(t0)}}
\[
  x(t) = \underbrace{x(t_0)}_{x_0} + \underbrace{v_m(t_0, t)}_{\frac{v(t) + v(t_0)}{2}}\overbrace{(t - t_0)}^{\Delta t}
  = x_0 + \frac{v(t) + v(t_0)}{2}\,(t - t_0)
\]
\[
  v(t) = \underbrace{v(t_0)}_{v_0} + \underbrace{a_m(t_0, t)}_{= a_c}\underbrace{(t - t_0)}_{\Delta t}
  = v_0 + a_c(t - t_0)
\]

\begin{dimostrazione}
(\textsuperscript{*} dalla nota adesiva:)
\[
  v_m(t, t_0) = \frac{v(t) + v(t_0)}{2} = \frac{v(t) + v_0}{2}
\]
\[
  v(t) = v_0 + a_c(t - t_0) \ \implies\ \overset{\textcolor{black!55}{v_m(t, t_0)}}{v_m}
  = \frac{v_0 + a_c(t - t_0) + v_0}{2}
  = v_0 + \frac12 a_c (t - t_0)
\]
\[
  x(t) = x_0 + v(t, t_0)(t - t_0) = x_0 +
  \underbrace{v_0(t - t_0) + \frac12 a_c (t - t_0)^2}_{\substack{\textcolor{black!55}{v_m(t, t_0)\cdot(t - t_0)} \\[0.6ex]
  \textcolor{black!55}{\left[v_0 + \frac12 a_c (t - t_0)\right](t - t_0)}}}
\]
\begin{formula}
x(t) = x_0 + v_0(t - t_0) + \frac12 a_c(t - t_0)^2
\end{formula}
\end{dimostrazione}
\begin{significato}
La posizione ha tre contributi: dove si parte ($x_0$), quanto si percorrerebbe
mantenendo la velocità iniziale ($v_0 \Delta t$) e il contributo dovuto all'accelerazione
($\frac12 a_c \Delta t^2$), che cresce col quadrato del tempo: per questo il grafico $x(t)$ è una parabola.
\end{significato}

\subsection{Relazione senza il tempo}
Una relazione importante si ha eliminando il tempo:
\[
  (t - t_0) = \frac{v(t) - v_0}{a_c}
  \ \Rightarrow\
  x(t) = x_0 + v_0\,\frac{v(t) - v_0}{a_c} + \frac12 a_c\left(\frac{v(t) - v_0}{a_c}\right)^2
\]
\begin{align*}
  x(t) &= x_0 + \frac{v_0\,(v(t) - v_0)}{a_c} + \frac12 \cancel{a_c}\,\frac{(v(t) - v_0)^2}{a_c^{\cancel{2}}} \\
       &= x_0 + \frac{v(t) - v_0}{a_c}\cdot\left[v_0 + \frac{v(t) - v_0}{2}\right]
       \qquad \left(\frac{2v_0 + v(t) - v_0}{2} = \frac{v(t) + v_0}{2}\right)
\end{align*}
\[
  x(t) = x_0 + \frac{v(t) - v_0}{a_c}\cdot\frac{v(t) + v_0}{2} = x_0 + \frac{v^2(t) - v_0^2}{2a_c}
  \ \Rightarrow\ x(t) - x_0 = \frac{v^2(t) - v_0^2}{2a_c}
\]
\begin{formula}
v^2(t) - v_0^2 = 2a_c\,(x(t) - x_0)
\end{formula}
\begin{significato}
Lega direttamente velocità e spazio percorso senza passare dal tempo:
utile quando si conosce quanto spazio c'è a disposizione (es.\ una frenata)
e si vuole sapere che velocità si raggiunge, o viceversa.
\end{significato}

% ---------------------------------------------------------------------
\section{Esempi: gravi}
\begin{esempio}{1 -- grave in caduta libera}
Legge moto: \ $y(t) = y_0 - \dfrac{g}{2}(t - t_0)^2$ \qquad (partenza da $y_0$ a $t_0$)

Assunzione:
\begin{itemize}[nosep]
  \item $v_0 = 0$
  \item Si trascura la resistenza dell'aria
  \item Accelerazione costante ($-g$)
\end{itemize}
\[
  v(t) = \frac{d}{dt}\left(y_0 - \frac{g}{2}(t - t_0)^2\right) = 0 - \frac{g}{\cancel{2}}\,\cancel{2}(t - t_0) = -g(t - t_0)
\]
\[
  a(t) = \frac{d}{dt}\left[-g(t - t_0)\right] = -g \quad \text{uniform.\ accelerato}
\]
% disegno: p15-y (appunti p. 15)
\begin{disegno}
\begin{tikzpicture}[disegno,x=0.9cm,y=0.9cm]
  \draw[asse] (0,0) -- (3.5,0) node[right] {$t$};
  \draw[asse] (0,0) -- (0,3.0) node[above right] {$y(t)$};
  \draw[tratt,black] (0,1.9) node[left] {$y_0$} -- (1.01,1.9) -- (1.01,0) node[below] {$t_0$};
  \draw[thick,black] (1.01,1.9) .. controls (1.4,1.75) and (1.82,1.1) .. (1.86,0);
  \node[punto,inner sep=1.1pt] at (1.01,1.9) {};
  \node[below,xshift=2pt] at (1.86,0) {$t_f$};
\end{tikzpicture}
\end{disegno}
% disegno: p15-va (appunti p. 15)
\begin{disegno}
\begin{tikzpicture}[disegno]
% --- v(t)
\node[anchor=east] at (-0.6,1.0) {$v(t)$:};
\draw[asse] (-0.4,0) -- (3.0,0) node[below right=1pt and 4pt] {$t$};
\draw[asse] (0,-0.9) -- (0,1.9);
\draw[curva,black] (1.0,0) -- (2.6,-0.75);
\draw[tratt] (2.6,-0.75) -- (2.6,0);
\node[above] at (1.0,0) {$t_0$};
\node[above] at (2.6,0) {$t_f$};
% --- a(t)
\begin{scope}[xshift=7.6cm]
\node[anchor=east] at (-1.1,1.1) {$a(t)$:};
\draw[asse] (-0.3,0) -- (4.2,0) node[right] {$t$};
\draw[asse] (0,-0.4) -- (0,2.0) node[right] {$a(t)$};
\draw[tratt] (0.9,0) -- (0.9,-0.75);
\draw[tratt] (3.6,0) -- (3.6,-0.75);
\draw[thick,black] (1.0,-0.95) -- (3.4,-0.95);
\node[above] at (1.2,0.1) {$t_0$};
\node[above] at (3.6,0.1) {$t_f$};
\end{scope}
\end{tikzpicture}
\end{disegno}
\textbf{NB!} Il grafico di $y(t)$ non è la traiettoria, ma come cambia il valore della funzione $y(t)$.

\medskip
$t_f = ?$ \qquad $y(t_f) = 0$ \implica $y_0 - \dfrac{g}{2}(t_f - t_0)^2 = 0$
\[
  y_0 = \frac{g}{2}(t_f - t_0)^2 \ \to\ \sqrt{\frac{2y_0}{g}} = \sqrt{(t_f - t_0)^2}
\]
\[
  |t_f - t_0| = \sqrt{\frac{2y_0}{g}} \qquad \text{assumiamo } t_f > t_0 \text{ sempre}
\]
\[
  t_f = t_0\,\sqrt{\frac{2y_0}{g}}
  \ \to\ v(t_f) = \underbrace{v_0}_{0} - g(t_f - t_0) = -g\left(t_0 + \sqrt{\frac{2y_0}{g}} - t_0\right)
  = -g\sqrt{\frac{2y_0}{g}} = -\sqrt{2y_0 g}
\]
\end{esempio}

\begin{esempio}{2 -- grave lanciato verso l'alto}
Grave lanciato verso l'alto con vel.\ $v_0$ e $a_c = -g$.
\[
  y(t) = y_0 + v_0(t - t_0) - \underbrace{\frac{g}{2}(t - t_0)^2}_{\text{termine quadr.\ } \to \text{ parabola}}
  \qquad
  v(t) = v_0 - g(t - t_0)
\]
Salendo $v(t) > 0$, scendendo $v(t) < 0$. In che tempo?
\begin{enumerate}[label=(\roman*)]
  \item $v(t) > 0 \to v_0 - g(t - t_0) > 0 \qquad v_0 > g(t - t_0)$
  \[
    v_0 > gt - gt_0 \ \to\ \frac{gt}{g} < \frac{v_0 + gt_0}{g}
    \qquad
    \underbrace{t < t_0 + \frac{v_0}{g}}_{\text{salita}} \ \to\ \text{apice}
  \]
  \item $v(t) < 0 \to \underbrace{t > t_0 + \dfrac{v_0}{g}}_{\text{discesa}}$
\end{enumerate}
% disegno: p16-grafici (appunti p. 16)
\begin{disegno}
\begin{tikzpicture}[disegno]
% --- x(t)
\draw[asse] (-0.6,0) -- (3.6,0);
\draw[asse] (0,-0.5) -- (0,2.8);
\draw[curva] (0.45,0) parabola bend (1.7,2.0) (2.95,0);
\fill[materia] (0.45,0) circle (1.3pt) (2.95,0) circle (1.3pt);
\draw[tratt] (0,2.0) -- (1.7,2.0);
\draw[tratt,->] (1.7,2.0) -- (1.7,-1.0);
\node[left] at (0,2.0) {apice};
\node[below] at (0.45,-0.05) {$t_0$};
\node[below] at (2.95,-0.05) {$t_f$};
\node[below] at (1.7,-1.0) {$t_0+\dfrac{v_0}{g}$};
\node[materia] at (1.7,3.1) {$x(t)$};
% --- v(t)
\begin{scope}[xshift=5.4cm]
\draw[asse] (-0.4,0.9) -- (3.8,0.9);
\draw[asse] (0,-0.5) -- (0,2.8);
\draw[thick,nota] (0.5,1.9) -- (2.4,-0.1);
\fill[nota] (0.5,1.9) circle (1.3pt) (2.4,-0.1) circle (1.3pt);
\draw[tratt] (0,1.9) -- (0.5,1.9);
\draw[tratt] (0,-0.1) -- (2.4,-0.1) -- (2.4,0.9);
\node[left] at (0,1.9) {$v_0$};
\node[left] at (0,-0.1) {$v_f$};
\node[below] at (0.5,0.85) {$t_0$};
\node at (2.9,1.5) {$t_f=t_0+\dfrac{2v_0}{g}$};
\node[nota] at (1.7,3.1) {$v(t)$};
\end{scope}
% --- a(t)
\begin{scope}[xshift=10.6cm]
\draw[asse] (-0.4,0.9) -- (3.6,0.9);
\draw[asse] (0,-0.5) -- (0,2.8);
\draw[thick,esempio] (0.5,0.4) -- (2.9,0.4);
\draw[tratt] (0,0.4) -- (0.5,0.4) -- (0.5,0.9);
\draw[tratt] (2.9,0.4) -- (2.9,0.9);
\node[left] at (0,0.4) {$-g$};
\node[above] at (0.5,0.9) {$t_0$};
\node[above] at (2.9,0.9) {$t_f$};
\node[esempio] at (1.7,3.1) {$a(t)$};
\end{scope}
\end{tikzpicture}
\end{disegno}
Apice a $t_0 + \dfrac{v_0}{g}$; \quad $t_f = t_0 + \dfrac{2v_0}{g}$.
\end{esempio}

\begin{delucidazione}
!!! \ Fine
\end{delucidazione}

% ---------------------------------------------------------------------
\section{Considerazioni}
\begin{enumerate}[label=(\Roman*)]
  \item $v$ e $a$ definite da coordinate cartesiane, ma la stessa definizione si può
        applicare a coordinate qualsiasi.\\
        Es.\ \implica \evid{ascissa curvilinea: $s(t)$}
  \[
    \begin{cases}
      v_s(t) = \displaystyle\lim_{\Delta t \to 0} \frac{s(t + \Delta t) - s(t)}{\Delta t}
             = \left.\frac{d}{dt}\,s(t)\right|_t = \dot{s}(t) \\[2.5ex]
      a_s(t) = \displaystyle\lim_{\Delta t \to 0} \frac{v_s(t + \Delta t) - v_s(t)}{\Delta t}
             = \left.\frac{d}{dt}\,v_s(t)\right|_t = \dot{v}_s(t) = \frac{d^2}{dt^2}\,s(t) = \ddot{s}(t)
    \end{cases}
  \]
  \item $v$ e $a$ sono chiamate \evid{velocità scalare} e \evid{accelerazione scalare}.\\
        \evid{Stessa cosa per $v_s$ e $a_s$}, il cui segno dipende dal verso di percorrenza di $s(t)$.
\end{enumerate}

% ===== da p05.tex =====
% =====================================================================
%  CAPITOLO 5 -- MOTO RETTILINEO VARIO E FORMULE GENERALI  (appunti pp. 18-19)
% =====================================================================
\chapter{Moto rettilineo vario e formule generali}

\section{Moto rettilineo vario}
$a(t)$ è una funzione non banale \implica \evid{$a(t)$ non costante}.

\begin{esempio}{moto smorzato esponenzialmente}
\[
  \underbrace{x(t) = x_0\left(1 - e^{-t/\tau}\right)}_{\text{spostamento}}
  \qquad [x_0] = [L] \qquad [\tau] = [T]
\]
% disegno: p18-smorzato (appunti p. 18)
\begin{disegno}
\begin{tikzpicture}[disegno]
\draw[asse] (0,0) -- (3.0,0) node[right] {$t$};
\draw[asse] (0,0) -- (0,2.6) node[right] {$x(t)$};
\node[below left] at (0,0) {$0$};
\draw[tratt,black] (0,1.75) -- (3.3,1.75);
\draw[curva,black] plot[domain=0:3.3,samples=60,smooth] (\x,{1.75*(1-exp(-\x/1.1))});
\end{tikzpicture}
\end{disegno}
\[
  v(t) = \dot{x}(t) = \frac{d}{dt}\,x_0\left(1 - e^{-t/\tau}\right) = \frac{d}{dt}\left(x_0 - x_0 e^{-t/\tau}\right)
\]
\[
  \frac{d}{dt}\,x_0 = 0
  \qquad
  \frac{d}{dt}\left(x_0 e^{-t/\tau}\right) = 0 - x_0\,\frac{d}{dt}\,e^{-t/\tau}
\]
\[
  D\!\left[e^{t} \circ \left(-\tfrac{1}{\tau}\,t\right)\right] = e^{-t/\tau}\left(-\frac{1}{\tau}\right) = -\frac{1}{\tau}\,e^{-t/\tau}
\]
\[
  \Rightarrow\ \underbrace{v(t) = \frac{x_0}{\tau}\,e^{-t/\tau}}_{\text{velocità}}
\]
\[
  a(t) = \frac{d}{dt}\left[\frac{1}{\tau}\cdot x_0 e^{-t/\tau}\right]
       = \frac{x_0}{\tau}\,\frac{d}{dt}\,e^{-t/\tau}
       = \frac{x_0}{\tau}\cdot\frac{1}{\tau}\,e^{-t/\tau}
\]
\[
  \underbrace{a(t) = \frac{x_0}{\tau^2}\,e^{-t/\tau}}_{\text{accelerazione}}
  \quad
  \begin{cases}
    < 0 \to v \text{ decrescente} \\
    > 0 \to v \text{ crescente}
  \end{cases}
\]
\end{esempio}

% ---------------------------------------------------------------------
\lezione{25/9/26}
\section{Formule generali cinematica a una dimensione}
\[
  x(t), \qquad v(t) = \dot{x}(t) \qquad \text{e} \qquad a(t) = \dot{v}(t) = \ddot{x}(t)
\]
\evid{Integrando}:
\[
  v(t) = v(t_0) + \int_{t_0}^{t} a(t')\,dt'
  \qquad \text{e} \qquad
  x(t) = x(t_0) + \int_{t_0}^{t} v(t')\,dt'
\]
\begin{formula}
x(t) = x(t_0) + \int_{t_0}^{t} \Big[\underbrace{v(t_0) + \int_{t_0}^{t'} a(t'')\,dt''}_{v(t')}\Big]\,dt'
\end{formula}
\begin{significato}
Integrare è l'operazione inversa del derivare: sommando tutte le piccole variazioni di
velocità ($a\,dt$) a partire da quella iniziale si ottiene la velocità in ogni istante;
sommando tutti i piccoli spostamenti ($v\,dt$) a partire dalla posizione iniziale si ottiene la posizione.
Basta quindi conoscere $a(t)$, $x(t_0)$ e $v(t_0)$ per ricostruire tutto il moto.
\end{significato}

\begin{delucidazione}
Attenzione: nell'integrale definito: \
$\displaystyle\int_{t_0}^{t} dt = t - t_0$
\end{delucidazione}

\begin{esempio}{moto uniform.\ acc.\ \implica $a(t) = a_c$}
\[
  x(t) = x(t_0) + \int_{t_0}^{t}\Big[v(t_0) + \int_{t_0}^{t'} a(t'')\,dt''\Big]\,dt'
\]
\[
  v(t') + \int_{t_0}^{t'} \underbrace{a(t'')}_{a_c}\,dt'' = v(t') + a_c\int_{t_0}^{t'} dt'' = v(t') + a_c(t' - t_0)
\]
\begin{align*}
  \int_{t_0}^{t} v(t_0) + a_c(t' - t_0)\,dt'
  &= \int_{t_0}^{t} v(t_0)\,dt' + \int_{t_0}^{t} a_c(t' - t_0)\,dt' \\
  &= v(t_0)\int_{t_0}^{t} dt' + a_c\left[\int_{t_0}^{t} t'\,dt' - t_0\int_{t_0}^{t} dt'\right] \\
  &= v(t_0)(t - t_0) + a_c\Big[\underbrace{\tfrac12 t^2 - \tfrac12 t_0^2 - t_0(t - t_0)}_{(\ast)}\Big]
\end{align*}
\[
  (\ast)\quad \tfrac12 t^2 - \tfrac12 t_0^2 - t_0 t + t_0^2 = \tfrac12 t^2 - t_0 t + \tfrac12 t_0^2 = \tfrac12(t - t_0)^2
\]
\[
  = v(t_0)(t - t_0) + \tfrac12 a_c(t - t_0)^2
\]
\[
  \Rightarrow\ x(t) = x(t_0) + v(t_0)(t - t_0) + \tfrac12 a_c(t - t_0)^2
\]
\end{esempio}

% ===== da p06.tex =====
% =====================================================================
%  CAPITOLO 6 -- MOTO ARMONICO   (appunti pp. 20-22)
% =====================================================================
\chapter{Moto armonico}

\begin{definizione}{Moto armonico}
Moto descritto da una \evid{funzione armonica del tempo} (seno e coseno).
\end{definizione}
\begin{formula}
x(t) = A\sin(\omega t + \varphi)
\end{formula}
\begin{significato}
Il corpo oscilla avanti e indietro tra $+A$ e $-A$ ripetendo sempre lo stesso andamento:
$A$ dice quanto si allontana dal centro, $\omega$ quanto velocemente oscilla,
$\varphi$ da che punto dell'oscillazione parte all'istante $t = 0$.
\end{significato}
(Se aggiungo $+\pi$, la funzione cambia segno.)
\[
  [\varphi] = [\text{angolo}] = \text{numero}
  \qquad
  [\omega] = [T]^{-1} \ \to\ [\omega t] = \text{adimensionale}
\]

\begin{itemize}
  \item \evid{``$\omega$'' \implica pulsazione} $\neq$ moto circolare: misura la rapidità con cui $\varphi$ varia
  \item \evid{``$\varphi$'' \implica fase} \implica angolo iniziale
  \item \evid{``$A$'' \implica ampiezza}
  \item \evid{``$x(t)$'' \implica elongazione}
\end{itemize}

% disegno: p20-grafico (appunti p. 20)
\begin{disegno}
\begin{tikzpicture}[disegno,yscale=0.65]
\def\A{2.0}
\draw[asse] (-3.4,0) -- (4.6,0) node[below right] {$t$};
\draw[asse] (0,-3.3) -- (0,3.4) node[right] {$x(t)$};
\draw[accento,thick,dash pattern=on 6pt off 7pt] (0,\A) -- (4.3,\A);
\draw[accento,thick,dash pattern=on 6pt off 7pt] (0,-\A) -- (5.3,-\A);
\node[accento,left] at (-0.15,\A) {$A$};
\node[accento,left] at (-0.15,-\A) {$-A$};
\draw[curva] plot[domain=0.02:4.25,samples=120,smooth] (\x,{-\A*cos(deg(2*pi/2.3*(\x-0.4)))});
\draw[curva,densely dotted] plot[domain=4.25:4.45,samples=10,smooth] (\x,{-\A*cos(deg(2*pi/2.3*(\x-0.4)))});
\draw[nota,thick] (0.4,-\A-0.1) -- (0.4,-\A-0.6) -- (2.7,-\A-0.6) -- (2.7,-\A-0.1);
\node[nota,below] at (1.55,-\A-0.65) {$T$};
\end{tikzpicture}
\end{disegno}

\begin{formula}
\omega T = 2\pi \ \to\ T = \frac{2\pi}{\omega}
\end{formula}
\begin{significato}
Il periodo è il tempo di un'oscillazione completa: in quel tempo l'argomento del seno
avanza di un giro intero ($2\pi$). Più grande è $\omega$, più breve è il periodo.
\end{significato}

\begin{itemize}
  \item \evid{``$T$'' \implica periodo} \qquad $[T] = [T]$ \quad s
  \item \evid{``$f$'' \implica frequenza} (``$\nu$'') \implica
        $f = \dfrac{\omega}{2\pi} = \dfrac{1}{T}$ \qquad $[f] = [T]^{-1}$ \quad $\frac{1}{\mathrm{s}}$
\end{itemize}
\begin{significato}
La frequenza conta quante oscillazioni complete avvengono in un secondo: è l'inverso del periodo.
\end{significato}

\section{Oscillazione}
\evid{Oscillazione}: $A \to -A \to A$ \implica \evid{Oscillazione completa} nel tempo $T$.

\section{Velocità}
\evid{Velocità}
\[
  v(t) = \frac{d}{dt}\,x(t) = \dot{x}(t) = \frac{d}{dt}\left[A\sin(\omega t + \varphi)\right] = A\cdot\omega\cdot\cos(\omega t + \varphi)
\]
\begin{formula}
v(t) = A\omega\cos(\omega t + \varphi)
\end{formula}
\begin{significato}
La velocità oscilla anch'essa, ma ``sfasata'' di un quarto di periodo rispetto alla posizione:
è massima ($A\omega$) quando il corpo passa per il centro ed è nulla agli estremi, dove inverte il moto.
\end{significato}
($\omega$ equilibra le dimensioni.)

\section{Accelerazione}
\evid{Accelerazione}
\[
  a(t) = \dot{v}(t) = \frac{d}{dt}\,A\omega\cos(\omega t + \varphi) = A\omega\left(-\omega\sin(\omega t + \varphi)\right)
\]
\[
  a(t) = -\omega^2\,\underbrace{A\sin(\omega t + \varphi)}_{x(t)}
\]
\begin{formula}
a(t) = -\omega^2\,x(t)
\end{formula}
\begin{significato}
L'accelerazione è sempre proporzionale allo spostamento dal centro e diretta in verso opposto:
più il corpo si allontana, più viene ``richiamato'' indietro. È questo richiamo che produce l'oscillazione.
\end{significato}
($\omega^2$ equilibra le dimensioni.)

\begin{nota}
$a(t) = \ddot{x}(t)$ \implica
\[
  \ddot{x}(t) = -\omega^2 x(t) \qquad \text{equazione differenziale in } x(t)
\]
\evid{Per essere un moto armonico deve soddisfare un'equazione differenziale.}
\end{nota}

\section{Osservazioni}
\evid{Osservazioni}
\begin{enumerate}[label=(\Roman*)]
  \item \evid{Il moto armonico è un caso speciale di moto periodico}:
        \[
          x(t) : x(t + T) \qquad T: \text{periodo del moto}
        \]
        (periodico $\supset$ armonico)\\
        \evid{Cosa differenzia un moto armonico da uno periodico \implica l'eq.~differenziale.}
  \item \evid{Moto circolare $\Leftrightarrow$ moto armonico}
  \item \evid{Come si può riscrivere!}
        \[
          x(t) = \alpha\sin(\omega t) + \beta\cos(\omega t) \qquad [\alpha] = [\beta] = [L]
        \]
\end{enumerate}

\begin{dimostrazione}
\[
  x(t) = A\sin(\omega t + \varphi) = \underbrace{A\sin(\omega t)\cos\varphi}_{\alpha} + \underbrace{A\cos(\omega t)\sin(\varphi)}_{\beta}
\]
\[
  \begin{cases}
    \alpha = A\cos\varphi \\
    \beta = A\sin\varphi
  \end{cases}
  \qquad A = \sqrt{\alpha^2 + \beta^2} \qquad \tan\varphi = \frac{\beta}{\alpha}
\]
\end{dimostrazione}
\begin{formula}
x(t) = \alpha\sin(\omega t) + \beta\cos(\omega t)
\end{formula}
Ci permette di ricavare lo sfasamento.
\begin{significato}
È lo stesso moto scritto come somma di un seno e di un coseno con coefficienti fissi:
$\alpha$ e $\beta$ si ricavano direttamente dalle condizioni iniziali, senza dover risolvere
equazioni con funzioni trigonometriche dell'angolo $\varphi$.
\end{significato}

\begin{delucidazione}
NB! È molto utile quando si conosce $t_0$ e $x_0$ e $v_0$, in quanto ci risparmia
passaggi trigonometrici non lineari.
\end{delucidazione}

% ---------------------------------------------------------------------
\section{Commenti generali su \texorpdfstring{$v$ e $a$}{v e a}}
\evid{Commenti generali su $v$ e $a$}
\begin{enumerate}
  \item ``speed'' \implica modulo della velocità;\\
        ``velocity'' \implica vettore velocità.
  \item Velocità fondamentale \implica $c$ ($\simeq 3\cdot 10^{8}\,\frac{\mathrm{m}}{\mathrm{s}}$) nel vuoto
        \begin{itemize}[nosep]
          \item $v \ll c$ \implica meccanica classica
          \item $v \sim c$ \implica meccanica relativistica
        \end{itemize}
  \item Velocità del suono (a STP): in aria $= 340\,\frac{\mathrm{m}}{\mathrm{s}}$
  \item Accelerazione: dinamica \implica legame accelerazione $\Leftrightarrow$ forze (Newton)
  \item Accelerazione fondamentale \implica $g = 9{,}81\,\frac{\mathrm{m}}{\mathrm{s}^2}$
        \implica non uguale ovunque
\end{enumerate}

% ===== da r1.tex =====
% =====================================================================
%  RIEPILOGO 1 -- MOTI IN UNA DIMENSIONE   (pagina aggiunta, non negli appunti)
% =====================================================================
\riepilogo{Moti in una dimensione}

Tutto parte dalla legge oraria $x(t)$: derivando si ottengono velocità e accelerazione,
integrando si torna indietro. Il tipo di moto dipende da com'è fatta l'accelerazione.

\begin{mappa}
  \node[nodoevid] (x) {$x(t)$\\ \small posizione};
  \node[nodo, right=1.6cm of x] (v) {$v(t) = \dot{x}$\\ \small velocità};
  \node[nodo, right=1.6cm of v] (a) {$a(t) = \dot{v} = \ddot{x}$\\ \small accelerazione};
  \draw[freccia] (x) to[bend left=20] node[above, font=\footnotesize] {deriv.} (v);
  \draw[freccia] (v) to[bend left=20] node[above, font=\footnotesize] {deriv.} (a);
  \draw[freccia] (a) to[bend left=20] node[below, font=\footnotesize] {integro} (v);
  \draw[freccia] (v) to[bend left=20] node[below, font=\footnotesize] {integro} (x);
  \node[nodo, below=2.2cm of x, xshift=-0.6cm] (mru) {\textbf{MRU}\\ \small $a = 0$\\ \small $v$ costante};
  \node[nodo, right=0.9cm of mru] (mrua) {\textbf{MRUA}\\ \small $a$ costante\\ \small $v$ retta};
  \node[nodo, right=0.9cm of mrua] (vario) {\textbf{Vario}\\ \small $a(t)$ qualsiasi};
  \node[nodo, right=0.9cm of vario] (arm) {\textbf{Armonico}\\ \small $a = -\omega^2 x$};
  \node[nodoevid, above=0.6cm of mrua, xshift=1.6cm] (tipo) {Che accelerazione?};
  \draw[freccia] (a.south) |- (tipo.east);
  \foreach \n in {mru, mrua, vario, arm} \draw[freccia] (tipo) -- (\n.north);
\end{mappa}

\begin{formulario}
v_m = \frac{x(t_1) - x(t_0)}{t_1 - t_0} & Velocità media: quanto si è spostato in media per unità di tempo (pendenza della secante di $x(t)$). \\
v = \frac{dx}{dt} & Velocità istantanea: rapidità di cambiamento della posizione in quell'istante (pendenza della tangente). \\
a = \frac{dv}{dt} = \frac{d^2x}{dt^2} & Accelerazione: rapidità di cambiamento della velocità. \\
x(t) = x_0 + v_c(t - t_0) & MRU: spazi uguali in tempi uguali, grafico $x(t)$ rettilineo. \\
v(t) = v_0 + a_c(t - t_0) & MRUA: la velocità cresce (o cala) linearmente nel tempo. \\
x(t) = x_0 + v_0\Delta t + \tfrac12 a_c\Delta t^2 & MRUA: posizione iniziale $+$ tratto a velocità iniziale $+$ contributo dell'accelerazione (parabola). \\
v^2 - v_0^2 = 2a_c(x - x_0) & MRUA senza il tempo: lega velocità e spazio percorso. \\
y(t) = y_0 - \tfrac{g}{2}(t - t_0)^2 & Caduta libera da ferma: accelerazione costante $-g$. \\
x(t) = x(t_0) + \int_{t_0}^{t} v\,dt' & Formula generale: sommando i piccoli spostamenti si ricostruisce la posizione. \\
x(t) = A\sin(\omega t + \varphi) & Moto armonico: oscillazione tra $-A$ e $+A$. \\
T = \frac{2\pi}{\omega},\ \ f = \frac{1}{T} & Periodo (durata di un'oscillazione) e frequenza (oscillazioni al secondo). \\
\ddot{x} = -\omega^2 x & Equazione del moto armonico: accelerazione proporzionale e opposta allo spostamento. \\
\end{formulario}

% ===== da p07.tex =====
% =====================================================================
%  CAPITOLO 7 -- CINEMATICA VETTORIALE IN 3 DIMENSIONI   (appunti pp. 23-25)
% =====================================================================
\chapter{Velocità vettoriale -- 3 dimensioni}

3-dim \implica 3 leggi orarie \implica $(x(t), y(t), z(t))$.

\[
  \vec{R}(t) = (x(t), y(t), z(t)) \ \to\ \evid{vettore posizione}
\]
% disegno: p23-vettori (appunti p. 23)
\begin{disegno}
\begin{tikzpicture}[disegno,scale=1.35]
\coordinate (O) at (0,0);
\coordinate (R0) at (1.08,1.55);
\coordinate (R1) at (1.98,1.98);
\draw[asse,black] (O) -- (0,2.1) node[right=2pt] {$z$};
\draw[asse,black] (O) -- (1.75,0) node[below right] {$y$};
\draw[asse,black] (O) -- (-1.15,-1.05) node[below right] {$x$};
\node[left] at (-0.12,0.2) {$O$};
\fill[esempio!15] (O) -- (R0) -- (R1) -- cycle;
\draw[accento,thick] plot[smooth,tension=0.7] coordinates {(0.43,0.95) (1.08,1.55) (1.98,1.98) (2.43,2.05)};
\draw[vett,nota] (O) -- (R1);
\draw[vett,materia] (O) -- (R0);
\draw[vett,esempio] (R0) -- (R1);
\node[materia] at (0.4,1.42) {$\vec R_0$};
\node[esempio,above] at (1.45,2.0) {$\Delta\vec R$};
\node[nota] at (1.72,0.85) {$\vec R_1$};
\end{tikzpicture}
\end{disegno}
\begin{itemize}[nosep]
  \item[] $\vec{R}_0 \to$ vettore posizione $P_0$ a $t_0$
  \item[] $\vec{R}_1 \to$ vettore posizione $P_1$ a $t_1$
\end{itemize}

\evid{Vettore spostamento}:
\[
  \Delta\vec{R} = \vec{R}_1 - \vec{R}_0
  \qquad
  \Delta\vec{R}(t_0, t_1) = \vec{R}(t_1) - \vec{R}(t_0)
\]

\section{Velocità media e istantanea}
\begin{definizione}{Velocità media (vettoriale)}
\[
  \vec{V}_m(t, t_0) = \frac{\Delta\vec{R}(t, t_0)}{t - t_0}
\]
\end{definizione}
\begin{definizione}{Velocità istantanea (vettoriale)}
\[
  \vec{V}(t) = \frac{d}{dt}\,\vec{R}(t) = \dot{\vec{R}}(t)
  = \left(\frac{d}{dt}x(t), \frac{d}{dt}y(t), \frac{d}{dt}z(t)\right)
  = (v_x, v_y, v_z) = \evid{(\dot{x}, \dot{y}, \dot{z})}
\]
\end{definizione}
\begin{significato}
In tre dimensioni la velocità è un vettore: le sue tre componenti sono le velocità con cui
cambiano $x$, $y$ e $z$; il suo modulo dice quanto si va veloci, la sua direzione dove si sta andando.
\end{significato}

\begin{delucidazione}
Definiamo la vel.\ media $\vec{V}_m(t, t + \Delta t) \doteq \dfrac{\Delta\vec{R}}{\Delta t}$.\\
Non ci dà informazioni sulla velocità lungo la traiettoria.\\
Facciamo tendere $\Delta t$ a $0$:
\[
  \vec{V}(t) \doteq \lim_{\Delta t \to 0} \vec{V}_m(t, t + \Delta t)
  = \lim_{\Delta t \to 0} \underbrace{\frac{\vec{R}(t + \Delta t) - \vec{R}(t)}{\Delta t}}_{\text{rapporto increm.\ in } \vec{R}}
  = \frac{d}{dt}\,\vec{R}(t) = \dot{\vec{R}}(t)
\]
NB! Il limite di una funzione vettoriale è equivalente a quello di una funzione scalare:
i versori $\hat{x}, \hat{y}$ e $\hat{z}$ cartesiani sono costanti, e il limite è lineare sulle
componenti del vettore.
\end{delucidazione}

\evid{Usando i versori su assi coordinati} $\hat{x}, \hat{y}, \hat{z}$:
\[
  \begin{cases}
    \vec{R} = \hat{x}\,x + \hat{y}\,y + \hat{z}\,z \\
    \vec{V} = \hat{x}\,v_x + \hat{y}\,v_y + \hat{z}\,v_z
  \end{cases}
\]

% disegno: p23-tangente (appunti p. 23)
\begin{disegno}
\begin{tikzpicture}[disegno,scale=1.25]
\coordinate (O) at (0,0);
\path (0,0) .. controls (0.3,1.75) and (1.5,3.5) .. (4.4,3.55)
  coordinate[pos=0.48] (P0) coordinate[pos=0.58] (P2)
  coordinate[pos=0.77] (P3) coordinate[pos=0.92] (P1);
\draw[asse,black] (O) -- (0,4.7);
\draw[asse,black] (O) -- (5.4,0);
\draw[asse,black] (O) -- (-1.7,-1.15);
\fill[accento!12] (O) -- (P0) -- (P2) -- cycle;
\draw[curva] (0,0) .. controls (0.3,1.75) and (1.5,3.5) .. (4.4,3.55);
\node[materia,rotate=10,anchor=west] at (4.5,3.62) {traiettoria};
\foreach \P in {P0,P2,P3,P1} \draw[->,thick,black] (O) -- (\P);
\draw[->,very thick,accento] (P0) -- (P2);
\draw[->,thick,esempio] (P0) -- (P1);
\draw[->,thick,esempio] ($(P0)!0.55!(P1)$) -- ++(0.32,-0.62) coordinate (dR);
\node[esempio,below right=-2pt] at (dR) {$\Delta\vec R$};
\draw[->,thick,accento] ($(P0)!0.5!(P2)$) -- (2.63,0.92);
\node[accento,right] at (2.63,0.92) {$\Delta R''$};
\draw[thick,black] (3.22,1.55) -- (3.3,1.42) -- (3.4,1.58);
\node[above left] at (P0) {$P_0$};
\node[above=2pt] at (P2) {$P_1''$};
\node[above=2pt] at (P3) {$P_1'$};
\node[above=2pt] at (P1) {$P_1$};
\end{tikzpicture}
\end{disegno}
Ogni volta che diminuisco $\Delta\vec{R}$ mi avvicino alla tangente alla curva
\implica velocità istantanea: vettore $\Delta\vec{R}$ si avvicina a $\vec{V}(t)$.

\begin{teorema}{Direzione della velocità}
``\evid{La direzione della velocità istantanea coincide con la tangente alla traiettoria}''
\end{teorema}

\begin{nota}
Non come per $v(t)$ in 1 dim.: era la tang.\ al graf.\ dello spost., non alla traiettoria. Come mai?
\end{nota}

\section{Formule generali}
\evid{Formule generali}
\begin{formula}
\vec{R}(t) = \vec{R}(t_0) + \int_{t_0}^{t} \vec{V}(t')\,dt'
\qquad \text{e} \qquad
\vec{V}(t) = \dot{\vec{R}}(t) = \frac{d}{dt}\,\vec{R}(t)
\end{formula}
\begin{significato}
Sono le stesse formule del caso unidimensionale applicate a ciascuna delle tre componenti:
la posizione si ottiene sommando, a partire da quella iniziale, tutti i piccoli spostamenti $\vec{V}\,dt$.
\end{significato}
Integrale di un vettore \implica vettore con 3 componenti che sono integrali.

\evid{Moto rettilineo uniforme} \implica
$\vec{R}(t) = \vec{R}(t_0) + \underbrace{\vec{V}_0}_{\text{vettore costante}}(t - t_0)$

\section{Accelerazione vettoriale}
\evid{Accelerazione vettoriale}
\begin{definizione}{Accelerazione vettoriale media}
\[
  \vec{a}_m(t, t + \Delta t) = \frac{\vec{V}(t + \Delta t) - \vec{V}(t)}{\Delta t}
\]
\end{definizione}
$\Delta t$ tende a $0$:
\[
  \lim_{\Delta t \to 0} a_m(t, t + \Delta t) = \lim_{\Delta t \to 0} \frac{v(t + \Delta t) - v(t)}{\Delta t}
  = \frac{d}{dt}\,\vec{V}(t) = \dot{\vec{V}}(t)
\]
Dunque\dots
\begin{formula}
\vec{a}(t) = \dot{\vec{V}}(t) = \frac{d}{dt}\,\vec{V}(t) = \ddot{\vec{R}}(t) = \frac{d^2}{dt^2}\,\vec{R}(t)
\end{formula}
\begin{significato}
L'accelerazione vettoriale misura come cambia il vettore velocità: cambia sia se cambia
il modulo (si va più o meno veloci) sia se cambia solo la direzione (si curva).
\end{significato}
\[
  \vec{a}(t) = (\dot{v}_x, \dot{v}_y, \dot{v}_z) = \hat{x}\,\dot{v}_x + \hat{y}\,\dot{v}_y + \hat{z}\,\dot{v}_z
  = (\ddot{x}, \ddot{y}, \ddot{z}) = \hat{x}\,\ddot{x} + \hat{y}\,\ddot{y} + \hat{z}\,\ddot{z}
\]

\subsection{Integrando}
\evid{Integrando}: \evid{velocità} ist.\ \implica
\[
  \vec{V}(t) = \vec{V}(t_0) + \int_{t_0}^{t} \vec{a}(t')\,dt'
\]
(si applica a tutti i componenti di $\vec{a}(t')$).

\evid{MRUA} \implica
$\vec{R}(t) = \vec{R}(t_0) + \vec{V}_0(t - t_0) + \dfrac{\vec{a}_c}{2}(t - t_0)^2$ \quad ($\vec{a}_c$: vettore costante)

\implica \evid{Eliminiamo la relazione dal tempo $(t - t_0)$}.\\
NB! Utile quando non so il tempo impiegato.

\begin{dimostrazione}
Dimostrare: $\vec{V}^2(t) = \vec{V}^2(t_0) + 2\vec{a}_c\,(\vec{R}(t) - \vec{R}(t_0))$.

Dunque partiamo da: $\vec{V}(t) = \vec{V}(t_0) + \vec{a}_c(t - t_0)$
\begin{align*}
  \vec{V}^2(t) &= \big[\underbrace{\vec{V}(t_0)}_{\vec{V}_0} + \vec{a}_c(t - t_0)\big]^2
               = \vec{V}_0^2 + 2\vec{V}_0\,\vec{a}_c(t - t_0) + \vec{a}_c^{\,2}(t - t_0)^2 \\
               &= \vec{V}_0^2 + 2\vec{a}_c\,\underbrace{\left(\vec{V}_0(t - t_0) + \tfrac12\vec{a}_c(t - t_0)^2\right)}_{\vec{R}(t) - \vec{R}(t_0)}
\end{align*}
\end{dimostrazione}
\begin{formula}
\vec{V}^2(t) = \vec{V}_0^2 + 2\vec{a}_c\,(\vec{R}(t) - \vec{R}(t_0))
\end{formula}
\begin{significato}
È la versione vettoriale di $v^2 - v_0^2 = 2a\,\Delta x$: il quadrato della velocità cambia in
base a quanto ci si sposta nella direzione dell'accelerazione, senza bisogno di conoscere il tempo.
\end{significato}

\section{Osservazioni}
\evid{Osservazioni}
\begin{enumerate}[label=(\roman*)]
  \item
  \[
    \underbrace{\int_{t_0}^{t} \vec{V}(t')\,dt'}_{\Delta\vec{R}(t, t_0)}
    \qquad\qquad
    \underbrace{l_{\text{tra.}} = \int_{t}^{t_0} |\vec{V}(t')|\,dt'}_{\text{lunghezza della traiettoria}}
  \]
  % disegno: p25-lunghezza (appunti p. 25)
  \begin{disegno}
\begin{tikzpicture}[disegno,scale=0.9]
\draw[asse,black] (0,0) -- (0,2.2);
\draw[asse,black] (0,0) -- (2.9,-0.1);
\draw[asse,black] (0,0) -- (-1.6,-1.15);
\draw[name path=c,thick,black] plot[smooth,tension=0.8] coordinates {(0.8,0.78) (1.4,1.28) (2.2,1.32) (2.9,1.05) (3.5,1.2) (3.9,1.32)};
\foreach \x/\b in {1.3/0.05,1.65/-0.05,2.0/0.0,2.4/0.1,2.75/-0.1,3.1/0.05,3.5/0.25,3.8/0.0}{
  \path[name path=v] (\x,-1) -- (\x,3);
  \path[name intersections={of=c and v,by=q}];
  \draw[materia,thick,dash pattern=on 2.5pt off 2.5pt] (q) -- (\x,\b);
  \fill[materia] (q) circle (1.4pt);
}
\draw[materia,thick] (1.65,1.42) to[bend left=50] (2.0,1.46);
\draw[->,materia,thick] (2.0,1.75) -- (2.45,2.45);
\node[materia,above] at (2.75,2.4) {$|\vec v(t)|\,dt$};
\end{tikzpicture}
  \end{disegno}
  Dato che ogni velocità è la tang.\ a un punto della traiettoria \implica la sommatoria di
  questi piccoli pezzettini ($\Sigma$) in modulo è la lung.\ della traiettoria.

  Mentre $\displaystyle\int_{t}^{t_0} |\vec{a}(t)|\,dt$ non ha un significato e scopo.
\end{enumerate}

% ===== da p08.tex =====
% =====================================================================
%  CAPITOLO 8 -- ACCELERAZIONE TANGENZIALE E NORMALE   (appunti p. 26)
% =====================================================================
\chapter{Accelerazione tangenziale e normale}

% disegno: p26-scomp1 (appunti p. 26)
\begin{disegno}
\begin{tikzpicture}[disegno]
\draw[asse,black] (0,0) -- (0,4.2) node[right=2pt] {$z$};
\draw[asse,black] (0,0) -- (4.2,0) node[below right] {$y$};
\draw[asse,black] (0,0) -- (-0.9,-0.9) node[below right] {$x$};
\draw[curva] plot[smooth,tension=0.7] coordinates {(0.81,0.56) (1.23,1.53) (1.8,2.02) (2.5,2.16) (3.38,1.98)};
\coordinate (P0) at (1.23,1.53);
\coordinate (P1) at (2.5,2.16);
\coordinate (A) at (2.0,2.56);
\coordinate (B) at (4.0,2.38);
\draw[vett,nota,thick] (P0) -- (A);
\draw[vett,nota,thick] (P1) -- (B);
\fill[nota] (P0) circle (1.6pt) (P1) circle (1.6pt);
\draw[vett,accento,thick] (A) -- (B);
\node[below right] at (P0) {$P_0$};
\node[nota,above left] at (A) {$\vec v(P_0)$};
\node[accento,above] at (3.1,2.5) {$\Delta\vec v$};
\node[nota,below right] at (3.75,2.33) {$\vec v(P_1)$};
\end{tikzpicture}
\end{disegno}
% disegno: p26-scomp2 (appunti p. 26)
\begin{disegno}
\begin{tikzpicture}[disegno,scale=0.95]
\coordinate (P0) at (0,0);
\coordinate (P1) at (2.56,0.99);
\coordinate (Q) at (2.17,2.2);
\coordinate (T) at (5.49,1.6);
\coordinate (Pe) at (3.39,0.99);
\draw[curva] plot[smooth,tension=0.7] coordinates {(-0.76,-1.26) (P0) (1.2,0.75) (P1) (3.8,1.0) (5.4,0.8) (7.4,0.34)};
\draw[vett,nota,thick] (P0) -- (Q);
\draw[vett,nota,thick] (P1) -- (T);
\fill[nota] (P0) circle (1.8pt) (P1) circle (1.8pt);
\draw[vett,esempio,thick] (Q) -- (4.14,3.77);
\draw[vett,esempio,thick] (Q) -- (Pe);
\draw[vett,accento,thick] (Q) -- (T);
\draw[esempio,thick,dash pattern=on 3pt off 4pt] (Pe) -- (4.6,-0.9);
\draw[esempio,thick] (4.6,-0.9) circle (1.8pt);
\draw[->,thick,black] (1.34,4.3) -- (1.34,3.55);
\node[below right] at (P0) {$P_0$};
\node[left=1pt] at (2.62,0.82) {$P_1$};
\node[nota] at (0.95,1.75) {$\vec v_1$};
\node[esempio,right] at (4.2,3.8) {$\overrightarrow{\Delta v}_{\parallel}$};
\node[accento] at (4.35,2.4) {$\Delta\vec v$};
\node[esempio] at (2.95,0.3) {$\Delta\vec v_{\perp}$};
\node[nota] at (4.95,1.15) {$\vec v_2$};
\end{tikzpicture}
\end{disegno}
Scomponiamo la vel.\ in due vettori ortogonali.

\begin{nota}
$\Delta\vec{V}_\perp$ è diretto ed è sempre verso la concavità della traiettoria.
\end{nota}

In cui $\Delta\vec{V}$ si scompone in:
\[
  \Delta\vec{V} \to
  \begin{cases}
    \Delta\vec{V}_{\parallel} \parallel \vec{V}(P_0) \\
    \Delta\vec{V}_{\perp} \perp \vec{V}(P_0)
  \end{cases}
\]

Come sappiamo dalla relazione:
\[
  a_m = \frac{\Delta v}{\Delta t} \ \to\ \vec{a}_m \propto \Delta\vec{V}
\]
\[
  \begin{cases}
    \vec{a}_{\parallel} \propto \Delta\vec{V}_{\parallel} \ \to\ \vec{a}_{\parallel} = \evid{\vec{a}_t} \ \to\ \text{``tangenziale''} \\
    \vec{a}_{\perp} \propto \Delta\vec{V}_{\perp} \ \to\ \vec{a}_{\perp} = \vec{a}_n \ \to\ \text{``normale''}
  \end{cases}
\]
($\vec{a}_t$ ha la direzione di $\vec{V}$, tang.\ alla traiettoria.)

\section{Raggio e centro di curvatura}
% disegno: p26-curvatura (appunti p. 26)
\begin{disegno}
\begin{tikzpicture}[disegno,scale=0.75]
\coordinate (P0) at (1.08,{-1.83-0.122*(1.08-3.2)^2});
\coordinate (P1) at (2.05,{-1.83-0.122*(2.05-3.2)^2});
\coordinate (A) at (2.92,-1.38);
\coordinate (C) at (4.33,-5.92);
\draw[curva] plot[domain=-2.1:6.67,samples=80,smooth] (\x,{-1.83-0.122*(\x-3.2)^2});
% raggi
\draw[nota,thick,dash pattern=on 5pt off 5pt] (P0) -- (C);
\draw[nota,thick,dash pattern=on 5pt off 5pt] (P1) -- (C);
% velocita'
\draw[->,nota,very thick] (P0) -- (A);
\draw[nota,very thick] (P1) -- (3.78,-1.66);
\fill[nota] (P0) circle (2pt) (P1) circle (1.6pt);
% accelerazione e componenti
\draw[esempio,thick,dash pattern=on 4pt off 4pt,-{Triangle[open,length=2.2mm,width=2mm]}] (A) -- (5.54,-0.6);
\draw[esempio,thick,dash pattern=on 4pt off 4pt] (A) -- (C);
\draw[->,esempio,thick] (A) -- ($(A)!0.13!(C)$);
\draw[->,accento,very thick] (A) -- (3.95,-1.58);
\fill (C) circle (2.2pt);
% etichette
\node[above left] at (P0) {$P_0$};
\node at (2.15,-2.65) {$P_1$};
\node[accento] at (4.45,-1.35) {$\vec a$};
\node[esempio,right] at (5.6,-0.55) {$\vec a_t$};
\node[esempio] at (4.05,-2.6) {$\vec a_N$};
\node at (1.4,-4.6) {$\rho$};
\node[below right] at (C) {$C$};
\draw[->,thick,black] (5.9,-6.25) to[out=10,in=-120] (8.4,-3.9);
\end{tikzpicture}
\end{disegno}
\begin{itemize}
  \item \evid{$\rho$ $=$ raggio di curvatura} (``ro'')
  \item \evid{$C$ $=$ centro di curvatura} \implica si trova nella regione di concavità della traiettoria
\end{itemize}

\begin{definizione}{Centro di curvatura}
È il centro della circonferenza meglio approssimabile alla traiettoria in un intervallo
infinitesimale tra due punti.
\end{definizione}

Se si avvicina $P_1$ a $P_0$, anche le perpend.\ si avvicinano.\\
Dunque\dots\ $\Delta\vec{V}(P_0)$ diventa uguale all'ortogonale $\to \vec{V}_t$.\\
Dunque $\vec{a}_n$ è diretto verso il centro di curvatura.

\[
  a_n \to \text{accelerazione centripeta}
\]

In un moto non rettilineo:
\begin{itemize}[nosep]
  \item $a_t$ \implica può essere zero (se costante in modulo)
  \item $a_n \neq 0$ sempre \implica la vel.\ cambia direzione
\end{itemize}

\begin{delucidazione}
NB! Cosa \evid{rappresentano}!
\begin{itemize}[nosep]
  \item[\implica] $a_t$ ``tangenziale'': come cambia il \evid{modulo} della velocità
  \item[\implica] $a_n$ ``\evid{normale}'': come cambia la \evid{direzione} della velocità
\end{itemize}
\end{delucidazione}

% ===== da p09.tex =====
% =====================================================================
%  CAPITOLO 9 -- MOTO DEI GRAVI LANCIATI (PROIETTILI)   (appunti pp. 27-29)
% =====================================================================
\chapter{Moto dei gravi lanciati (proiettili)}

Moto bidimensionale.

% disegno: p27-lancio (appunti p. 27)
\begin{disegno}
\begin{tikzpicture}[disegno]
  \draw[asse] (-0.4,0) -- (3,0) node[below right] {$x$};
  \draw[asse] (0,-0.5) -- (0,2.2) node[above] {$y$};
  \draw[vett,black] (0,0) -- (1.6,1.0) node[above right] {$\vec{v}_0$};
  \draw[thick] (0.9,0) arc[start angle=0,end angle=32,radius=0.9];
  \node at (1.15,0.25) {$\theta$};
  \draw[vett,black] (3.7,2.0) -- (3.7,1.0);
  \node at (4.0,2.0) {$\vec{g}$};
\end{tikzpicture}
\end{disegno}
\[
  \vec{g} = -g\,\hat{y}
  \qquad
  \text{Legge oraria} =
  \begin{cases}
    x = x_0 + v_0\cos\theta\cdot t \\
    y = y_0 + v_0\sin\theta\cdot t - \frac12 g t^2
  \end{cases}
\]
$x_0 = y_0 = 0$ \implica
\begin{formula}
\begin{cases}
  x(t) = v_0\cos\theta\; t \\
  y(t) = v_0\sin\theta\; t - \underbrace{\tfrac12 g t^2}_{\text{termine parabolico con concavità verso il basso}}
\end{cases}
\end{formula}
\begin{significato}
Il moto si divide in due moti indipendenti: in orizzontale il proiettile va a velocità
costante ($v_0\cos\theta$, nessuna accelerazione), in verticale sale e ricade come un grave
lanciato verso l'alto con velocità $v_0\sin\theta$ e accelerazione $-g$.
\end{significato}

\section{Due fasi}
\evid{Due fasi}:
\begin{itemize}
  \item \evid{Salita}: $v_y(t) > 0 \to t < \dfrac{v_0\sin\theta}{g}$
        \implica $t_s = \dfrac{v_0\sin\theta}{g}$ per $v_y = 0$ (tempo di salita)
  \item \evid{Discesa}: $v_y(t) < 0 \to t > \dfrac{v_0\sin\theta}{g}$
        \implica $v_y(t) = \underbrace{-g\left(t - \dfrac{v_0\sin\theta}{g}\right)}_{\text{formulazione più ``elegante''}}$
\end{itemize}

\section{Quota massima}
\evid{Quota massima}
\[
  y(t = t_s) = (v_0\sin\theta)\overbrace{\left(\frac{v_0\sin\theta}{g}\right)}^{t} - \frac{g}{2}\overbrace{\left(\frac{v_0\sin\theta}{g}\right)^2}^{t^2}
  = \frac{v_0^2\sin^2\theta}{2g}
\]
Dunque:
\begin{formula}
h = \frac{v_0^2\sin^2\theta}{2g}
\end{formula}
\begin{significato}
L'altezza raggiunta dipende solo dalla componente verticale della velocità iniziale
($v_0\sin\theta$): più si lancia in verticale e più forte, più si sale; la parte orizzontale non conta.
\end{significato}
con vel.\ verticale iniziale \implica $v_y(t_s) = 0$.

Il moto è \evid{asimmetrico}:
\begin{itemize}[nosep]
  \item a $t_0$: \ $v_y \neq 0$ e $h = 0$
  \item a $t_s$: \ $v_y = 0$ e $h \neq 0$
\end{itemize}

\section{Tempo totale di volo}
\evid{Tempo totale di volo}

$y(t_0) = 0 \to y(t_{tot}) = 0$ \implica trovo le \evid{radici della legge oraria $y(t)$}.
\[
  y(\bar{t}) = v_0\sin\theta\,\bar{t} - \tfrac12 g\bar{t}^2 = 0
  \qquad
  \bar{t}\left(v_0\sin\theta - \tfrac12 g\bar{t}\right) = 0
\]
Per $\bar{t} = 0$ \implica punto di partenza.
\begin{formula}
\bar{t} = \frac{2v_0\sin\theta}{g} = t_{tot}
\end{formula}
\begin{significato}
Il proiettile resta in aria il doppio del tempo che impiega a salire: salita e discesa durano uguale.
\end{significato}
Doppio del tempo di salita, dunque:
\[
  t_s = t_d = \frac{v_0\sin\theta}{g} \qquad (t_s: \text{tempo di salita})
\]
\begin{delucidazione}
La vel.\ finale su $y$ è $v_y = v_0\sin\theta$.
\end{delucidazione}

\section{Gittata}
\evid{Gittata}
\[
  x(t_{tot}) = v_0\cos\theta\cdot t_{tot} = v_0\cos\theta\cdot\frac{2v_0\sin\theta}{g}
  = \frac{v_0^2\cdot\overbrace{2\cos\theta\sin\theta}^{\sin(2\theta)}}{g}
\]
\begin{formula}
x(t_{tot}) = \frac{v_0^2\sin(2\theta)}{g}
\end{formula}
\begin{significato}
La distanza orizzontale percorsa dipende dall'angolo tramite $\sin(2\theta)$: è massima
a $45^\circ$ e uguale per due angoli che sommati danno $90^\circ$.
\end{significato}

\begin{nota}
$\sin(2\theta) = \sin\!\left(2\left(\theta + \frac{\pi}{2}\right)\right) = \sin(2\theta)$
% disegno: p28-angoli (appunti p. 28)
\begin{disegno}
\begin{tikzpicture}[disegno,black!55,thick]
  \draw (0,0) ellipse[x radius=1.75,y radius=1.45];
  \draw (-1.95,0.05) -- (2.5,-0.08);
  \draw (0,-1.75) -- (0,1.85);
  \draw (0,0) -- (30:1.6);
  \draw (0,0) -- (148:1.55);
  \draw[dashed] (175:0.85) arc[start angle=175,end angle=32,radius=0.85];
  \draw[dashed] (125:0.55) arc[start angle=125,end angle=45,radius=0.55];
  \draw (0.9,0) arc[start angle=0,end angle=30,radius=0.9];
  \node at (13:1.18) {$\theta$};
\end{tikzpicture}
\end{disegno}
Lanciare un proiettile con $\theta = 30^\circ$ e $\theta = 60^\circ$: è uguale la gittata.
\end{nota}

\section{Traiettoria}
\evid{Traiettoria}
\[
  \begin{cases}
    x = v_0\cos\theta\cdot t \ \longmapsto\ t = \dfrac{x}{v_0\cos\theta} \\[1.5ex]
    y = v_0\sin\theta\cdot t - \dfrac{g}{2}t^2 \ \longrightarrow\
        v_0\sin\theta\cdot\dfrac{x}{v_0\cos\theta} - \dfrac{g}{2}\,\dfrac{x^2}{v_0^2\cos^2\theta}
  \end{cases}
\]
\[
  \frac{\sin\theta}{\cos\theta} = \tan\theta
  \ \Rightarrow\
  y = x\cdot\tan\theta - \frac{g}{2}\left(\frac{x}{v_0\cos\theta}\right)^2
\]
\begin{formula}
y(x) = x\tan\theta - \frac12 g\,\frac{x^2}{v_0^2\cos^2\theta}
\end{formula}
\begin{significato}
È l'equazione della traiettoria vera e propria (senza il tempo): una parabola rivolta verso il basso.
Il primo termine è la retta lungo cui si partirebbe senza gravità, il secondo è quanto la gravità
``piega'' il percorso verso il basso, sempre di più man mano che ci si allontana.
\end{significato}
Parabola (termine quadratico).

% disegno: p28-traiettoria (appunti p. 28)
\begin{disegno}
\begin{tikzpicture}[disegno,
    tang/.style={->,thick,nota,>={Stealth[length=1.6mm]}}]
  % parabola: y = H(1-((x-a)/a)^2), a=2, H=3
  \def\a{2}\def\H{3}
  \draw[thick,black,-{Latex[open,length=2.2mm]}] (-0.6,0) -- (5.6,0) node[right] {$x$};
  \draw[thick,black,-{Latex[open,length=2.2mm]}] (0,-0.45) -- (0,4.2) node[above right] {$y$};
  \draw[dashed,black] (0,\H) -- (1.5,\H);
  \node[left] at (0,\H) {$y(t_s)$};
  \draw[curva,very thick,materia!70] plot[domain=0:4,samples=60] (\x,{\H*(1-((\x-\a)/\a)^2)});
  % frecce tangenti (gialle)
  \foreach \xx in {0.2,0.65,1.15,1.6,2.55,3.15,3.65}{
    \pgfmathsetmacro{\yy}{\H*(1-((\xx-\a)/\a)^2)}
    \pgfmathsetmacro{\m}{-2*\H*(\xx-\a)/(\a*\a)}
    \pgfmathsetmacro{\L}{0.55/sqrt(1+\m*\m)}
    \draw[tang] (\xx,\yy) -- ++(\L,{\m*\L});
  }
  \draw[tang] (\a,\H) -- (3.75,\H-0.08) node[above right,nota] {$\vec{v}_s$};
  \fill[nota] (\a,\H) circle (1.6pt);
  \fill[black] (4,0) circle (1.6pt);
  \node[below] at (4.15,-0.05) {$G=x(t_{\mathrm{tot}})$};
\end{tikzpicture}
\end{disegno}

\subsection{Calcoliamo la gittata}
\evid{Calcoliamo la gittata}
\[
  y(x = G) = 0 = x\left[\tan\theta - \frac12 g\,\frac{x}{v_0^2\cos^2\theta}\right]
  \qquad (x = 0: \text{caso iniziale a } t = 0)
\]
\[
  \tan\theta - \frac{g}{2}\,\frac{x}{v_0^2\cos^2\theta} = 0
  \ \to\
  x = \tan\theta\cdot\frac{2}{g}\,v_0^2\cos^2\theta
  = 2\,\frac{\sin\theta}{\cancel{\cos\theta}}\cdot\cos^{\cancel{2}}\theta\cdot\frac{v_0^2}{g}
  = \frac{v_0^2}{g}\,\sin(2\theta) \ \checkmark
\]
($2\sin\theta\cos\theta = \sin(2\theta)$)

\subsection{Altezza massima}
\evid{Altezza massima}

$y_{max} = y(t_s)$ oppure $y'(x) = 0$. Dove? In
$x_v = \dfrac{G}{2} = \dfrac{v_0^2\sin(2\theta)}{2g}$.
\begin{align*}
  \left.\frac{d}{dx}\,y(x)\right|_{x = x_v}
  &= \frac{d}{dx}\left[x\tan\theta - \frac{g}{2}\left(\frac{x}{v_0\cos\theta}\right)^2\right]
   = \tan\theta\,\frac{d}{dx}x - \frac{g}{2v_0^2\cos^2\theta}\,\frac{d}{dx}x^2 \\
  &= \tan\theta - \frac{g}{2v_0^2\cos^2\theta}\cdot 2x = \tan\theta - \frac{gx}{v_0^2\cos^2\theta} = 0
\end{align*}
\[
  x_v = \tan\theta\,\frac{v_0^2\cos^2\theta}{g} = \frac{v_0^2\sin(2\theta)}{2g} = \frac{G}{2}
  \qquad (\text{formula di duplicazione})
\]
\begin{align*}
  y_{max} = y(x_v)
  &= \frac{v_0^2\sin(2\theta)}{2g}\,\tan\theta - \frac{g}{2}\,\frac{v_0^4\sin^2(2\theta)}{4g^2v_0^2\cos^2\theta}
   \qquad (\sin^2(2\theta) = 4\cos^2\theta\sin^2\theta) \\
  &= \frac{v_0^2\,2\sin\theta\cos\theta}{2g}\,\frac{\sin\theta}{\cos\theta}
     - \frac{g}{2}\,\frac{v_0^2}{4g^2}\,\frac{4\cos^2\theta\sin^2\theta}{\cos^2\theta} \\
  &= \frac{v_0^2\sin^2\theta}{g} - \frac{v_0^2\sin^2\theta}{2g} = \frac{v_0^2\sin^2\theta}{2g} = y_{max}
\end{align*}

\begin{delucidazione}
\textbf{Valori massimi}: al variare di $\theta$, possiamo avere $h_{max}$ con
$\theta = \frac{\pi}{2}$ ($= 90^\circ$) \implica $h = \dfrac{v_0^2}{g}$.\\
Invece la \evid{gittata massima} si ha a $\theta = \frac{\pi}{4} \approx 45^\circ$.
\end{delucidazione}

\section{Ellisse dei vertici}
\evid{Ellisse dei vertici}

Cerchiamo il \evid{luogo geometrico dei vertici della parabola al variare di $\theta$}.
\[
  \begin{cases}
    x_v = \dfrac{G}{2} = \dfrac{v_0^2\sin\theta\cos\theta}{g} \leq \dfrac{v_0^2}{2g} & \text{(massimo)} \\[2ex]
    2y_v = \dfrac{v_0^2\sin^2\theta}{g} \leq \dfrac{v_0^2}{2g} & \text{(massimo)}
  \end{cases}
\]
Liberiamoci di $\theta$:
\begin{align*}
  (x_v)^2 + (2y_v)^2
  &= \frac{v_0^4\sin^2\theta\cos^2\theta}{g^2} + \frac{v_0^4\sin^4\theta}{g^2}
   = \frac{v_0^4\sin^2\theta\cos^2\theta + v_0^4\sin^4\theta}{g^2} \\
  &= \frac{v_0^4\sin^2\theta}{g}\,\underbrace{(\cos^2\theta + \sin^2\theta)}_{= 1}
   = \underbrace{\frac{v_0^2\sin^2\theta}{g}}_{= 2y_v}\,\frac{v_0^2}{g} = 2y_v\,\frac{v_0^2}{g}
\end{align*}
\begin{formula}
x_v^2 + 4y_v^2 - 2y_v\,\frac{v_0^2}{g} = 0 \qquad \text{(ellisse)}
\end{formula}
\begin{significato}
Cambiando l'angolo di lancio (a parità di $v_0$) i punti più alti delle traiettorie
non stanno in posti qualsiasi, ma tutti su questa ellisse.
\end{significato}

``Viene fuori'' un'ellisse, che:
\begin{itemize}[nosep]
  \item[\implica] passa per l'origine
  \item[\implica] ha centro $C\left(0;\ \dfrac{v_0^2}{4g}\right)$ (metà altezza massima); massimo $y_v = \dfrac{v_0^2}{2g}$
\end{itemize}
% disegno: p29-ellisse (appunti p. 29)
\begin{disegno}
\begin{tikzpicture}[disegno,x=0.12mm,y=-0.12mm]
  % assi
  \draw[asse,black] (100,275) -- (465,275) node[below right] {$x_v$};
  \draw[asse,black] (262,300) -- (262,95) node[right] {$y_v$};
  % ellisse (blu)
  \draw[very thick,materia!75]
     (262,135) .. controls (180,145) and (125,175) .. (125,205)
               .. controls (125,235) and (180,260) .. (262,270)
               .. controls (320,255) and (365,235) .. (385,210)
               .. controls (365,185) and (320,155) .. cycle;
  % retta tratteggiata verticale con doppio segno
  \draw[dashed,black] (325,160) -- (325,355);
  \draw[black] (319,225) -- (319,300);
  \draw[black] (331,225) -- (331,300);
  \draw[black,thick] (316,358) -- (334,358) -- (325,378) -- cycle;
  \fill (325,160) circle (1.4pt);
  % tratteggio dal vertice destro
  \draw[dashed,black] (385,212) -- (385,275);
  % freccia gialla
  \draw[->,very thick,nota] (262,205) .. controls (310,195) and (345,200) .. (385,212)
        .. controls (420,225) and (460,225) .. (500,215);
  \fill[nota] (262,205) circle (2.2pt);
  \node[nota,left] at (255,190) {$c$};
  \fill (385,212) circle (1.7pt);
  % freccia nera verso v0^2/2g
  \draw[->,thick,black] (385,212) .. controls (400,190) and (415,170) .. (442,155);
  \node at (472,130) {$\dfrac{v_0^2}{2g}$};
  \draw[black!55,thick] (510,95) .. controls (522,98) and (518,118) .. (521,128)
        .. controls (518,138) and (522,158) .. (510,162);
  \node[black!55,rotate=8,anchor=west] at (528,108) {massimo};
  % punto (0; v0^2/4g)
  \node[nota,anchor=west] at (525,195) {\large$\left(0;\,\dfrac{v_0^2}{4g}\right)$};
  \draw[black!55,thick] (705,160) .. controls (717,163) and (712,182) .. (715,195)
        .. controls (712,208) and (717,235) .. (697,240);
  \node[black!55,anchor=west] at (730,195) {metà altezza};
  \node[black!55,anchor=west] at (790,235) {massima};
\end{tikzpicture}
\end{disegno}
Se pongo $x$, incontra l'ellisse in due punti: sono i due angoli complementari (es.\ $30^\circ$ e $60^\circ$).

% ===== da p10.tex =====
% ---------------------------------------------------------------------
%  (segue CAPITOLO 9)  Parabola di sicurezza e lancio orizzontale (appunti pp. 30-31)
% ---------------------------------------------------------------------
\section{Parabola di sicurezza}
\evid{Parabola di sicurezza}

Si cerca di \evid{determinare i punti dello spazio non raggiungibili con $v_0$ e $\theta$}.
\[
  \text{traiettoria}\quad y = x\tan\theta - \frac{g}{2}\,\frac{x^2}{v_0^2\cos^2\theta}
  \qquad \left(\to \frac{1}{\cos^2\theta}\right)
\]
\[
  \frac{1}{\cos^2\theta} = \frac{\sin^2\theta + \cos^2\theta}{\cos^2\theta}
  = \underbrace{\frac{\cos^2\theta}{\cos^2\theta}}_{= 1} + \underbrace{\frac{\sin^2\theta}{\cos^2\theta}}_{= \tan^2\theta}
  = 1 + \tan^2\theta
\]
\[
  \Rightarrow\ y = x\tan\theta - \frac{g}{2}\,\frac{x^2}{v_0^2}\,(1 + \tan^2\theta)
\]
Ciò l'abbiamo fatto per avere un'equazione sol.\ in $\tan\theta$, per poi risolvere per quella.

\begin{delucidazione}
Si ottiene un'equazione quadratica in $\tan\theta$.
\end{delucidazione}

\[
  x\tan\theta - \frac{gx^2}{2v_0^2} - \frac{gx^2}{2v_0^2}\tan^2\theta - y = 0
\]
\[
  \Rightarrow\ \underbrace{\frac{gx^2}{2v_0^2}}_{a}\tan^2\theta \underbrace{- x}_{b}\tan\theta
  + \underbrace{\left(\frac{gx^2}{2v_0^2} + y\right)}_{c} = 0
  \qquad \text{(ho invertito i segni)}
\]
\begin{align*}
  \tan\theta
  &= \frac{x \pm \sqrt{x^2 - \cancel{4}^{\,2}\,\dfrac{gx^2}{\cancel{2}v_0^2}\left(\dfrac{gx^2}{2v_0^2} + y\right)}}{\cancel{2}\,\dfrac{gx^2}{\cancel{2}v_0^2}}
   = \frac{v_0^2}{gx^2}\left(x \pm \sqrt{x^2 - \frac{2x^2}{v_0^2}\left(\frac{gx^2}{2v_0^2} + y\right)}\right) \\
  &= \frac{v_0^2}{gx^2}\left(x \pm \sqrt{x^2\left[1 - \frac{2g}{v_0^2}\left(\frac{gx^2}{2v_0^2} + y\right)\right]}\right)
   = \frac{v_0^2}{gx^2}\left(x \pm x\sqrt{1 - \frac{2g}{v_0^2}\left(\frac{gx^2}{2v_0^2} + y\right)}\right)
   \quad (x > 0) \\
  &= \frac{v_0^2}{gx^{\cancel{2}}}\,\cancel{x}\left(1 \pm \sqrt{1 - \frac{2g}{v_0^2}\left(\frac{gx^2}{2v_0^2} + y\right)}\right)
   = \frac{v_0^2}{g}\Bigg(1 \pm \sqrt{\underbrace{1 - \frac{2g}{v_0^2}\left(\frac{gx^2}{2v_0^2} + y\right)}_{\Delta}}\,\Bigg)
\end{align*}

\evid{Cosa cerchiamo?} Quell'\evid{angolo per cui il grave non raggiunga mai un punto $(x, y)$
dato un certo $\theta$ e $v_0$, dunque vogliamo che l'eq.~della traiettoria non abbia radici}
\implica \evid{$\Delta < 0$}.
\[
  \begin{cases}
    \Delta > 0 \to 2 \text{ rad.} \\
    \Delta = 0 \to 1 \text{ rad.} \\
    \Delta < 0 \to 0 \text{ rad.}
  \end{cases}
\]
\[
  \Rightarrow\ 1 - \frac{2g}{v_0^2}\left(\frac{gx^2}{2v_0^2} + y\right) < 0
  \qquad \text{Risolviamo in } y
\]
\[
  \to\ 1 - \frac{2g^2x^2}{2v_0^4} + \frac{2gy}{v_0^2} < 0
  \qquad
  -\frac{2gy}{v_0^2} > 1 - \frac{2g^2x^2}{2v_0^4}
  \qquad
  \frac{2gy}{v_0^2} < \frac{2g^2x^2}{2v_0^4} - 1
\]
\[
  \to\ y < \frac{v_0^2}{2g}\,\frac{2g^2x^2}{2v_0^4} - \frac{v_0^2}{2g}
  \ \to\
  \underbrace{\evid{y > \dfrac{v_0^2}{2g} - \dfrac{gx^2}{2v_0^2}}}_{\text{zona di sicurezza}}
\]
\begin{formula}
y = \frac{v_0^2}{2g} - \frac{gx^2}{2v_0^2}
\end{formula}
\begin{significato}
Con una data velocità di lancio $v_0$, qualunque sia l'angolo, il proiettile non può superare
questa parabola: i punti sopra (zona di sicurezza) non vengono mai raggiunti, quelli sotto
possono essere colpiti con uno o due angoli diversi.
\end{significato}
% disegno: p30-zone (appunti p. 30)
\begin{disegno}
\begin{tikzpicture}[disegno,x=0.17mm,y=-0.17mm]
  % parabola: piedi a 350+-115, vertice (350,220)
  \def\parab{plot[domain=233:467,samples=50] (\x,{312-92*(1-((\x-350)/117)^2)})}
  % zona di sicurezza (verde): fuori dalla parabola
  \begin{scope}[even odd rule]
    \clip (205,312) -- (205,240) .. controls (230,170) and (290,140) .. (330,140)
          -- (500,140) -- (505,290) -- (480,312) -- cycle
          \parab -- cycle;
    \foreach \i in {0,...,16} \draw[esempio!80,thick] (190+\i*20,320) -- ++(45,-180);
  \end{scope}
  % zona colpibile (rossa): sotto la parabola
  \begin{scope}
    \clip \parab -- cycle;
    \foreach \i in {0,...,11} \draw[accento!85,thick] (230+\i*20,322) -- ++(25,-100);
  \end{scope}
  % assi
  \draw[thick,black,-{Latex[open,length=2.4mm]}] (178,312) -- (572,312);
  \draw[thick,black,-{Latex[open,length=2.4mm]}] (350,412) -- (350,103);
  % parabola
  \draw[very thick,materia!75] \parab;
  % quota massima
  \draw[dashed,materia!75,thick] (183,232) -- (350,220);
  \fill[materia!75] (183,232) circle (1.4pt);
  \node[materia!75,left] at (175,232) {$\dfrac{v_0^2}{2g}$};
  \node[materia!75] at (468,372) {$\dfrac{v_0^2}{g}$};
  % equazione
  \draw[->,thick,materia!75] (462,306) .. controls (490,290) and (510,282) .. (535,272);
  \node[materia!75,anchor=west] at (530,262) {$y=\dfrac{v_0^2}{2g}-\dfrac{g\,x^2}{2v_0^2}$};
  % etichette zone
  \draw[->,thick,esempio] (455,195) -- (478,163);
  \node[esempio,anchor=west] at (400,132) {zona di sicurezza};
  \draw[->,thick,accento] (292,306) -- (258,360);
  \node[accento,align=left,anchor=north west] at (148,375) {zona\\colpibile};
\end{tikzpicture}
\end{disegno}

\begin{delucidazione}
$\Delta < 0$: il punto non è raggiungibile \implica parabola di sicurezza.\\
All'interno della gittata, ma dove non vieni colpito? \implica Mettiamo $\Delta < 0$.
\end{delucidazione}

% ---------------------------------------------------------------------
\section{Grave lanciato orizzontalmente da una certa altezza \texorpdfstring{$h$}{h}}
\evid{Grave lanciato orizzontalmente da una certa $h$}
\[
  x_0 = 0 \qquad y_0 > 0 \qquad \theta = 0
\]
Legge oraria:
\[
  \begin{cases}
    x(t) = v_0 t \\
    y(t) = \underbrace{y_0}_{h} - \frac12 g t^2
  \end{cases}
\]
\begin{formula}
y(x) = h - \frac12 g\left(\frac{x}{v_0}\right)^2
\end{formula}
\begin{significato}
Il grave parte in orizzontale e cade sempre più ripido: in orizzontale avanza a velocità
costante $v_0$, in verticale cade come in caduta libera partendo da fermo.
\end{significato}
% disegno: p31-orizz (appunti p. 31)
\begin{disegno}
\begin{tikzpicture}[disegno,x=0.24mm,y=0.24mm,
    vb/.style={->,very thick,materia!75,>={Stealth[length=1.8mm]}}]
  \def\H{108}\def\L{178}
  \draw[asse,black] (-25,0) -- (220,0);
  \draw[asse,black] (0,-30) -- (0,160);
  \draw[thick,black] plot[domain=12:\L,samples=50] (\x,{\H*(1-(\x/\L)^2)});
  % velocita' tangenti
  \draw[vb] (12,{\H*(1-(12/\L)^2)}) -- ++(75,0) node[above,pos=0.85,yshift=10pt] {$\vec{v}_0$};
  \foreach \xx/\l in {105/50, 155/45}{
    \pgfmathsetmacro{\yy}{\H*(1-(\xx/\L)^2)}
    \pgfmathsetmacro{\m}{-2*\H/\L*\xx/\L}
    \pgfmathsetmacro{\d}{\l/sqrt(1+\m*\m)}
    \draw[vb] (\xx,\yy) -- ++(\d,{\m*\d});
  }
\end{tikzpicture}
\end{disegno}

$x$ finale?
\[
  y(x) = h - \frac12 g\left(\frac{x}{v_0}\right)^2
  \qquad
  h - \frac12 g\,\frac{x^2}{v_0^2} = 0
  \qquad
  \sqrt{x^2} = \sqrt{h\cdot 2\cdot\frac{v_0^2}{g}}
  \ \to\
  \evid{x = (\pm)\,v_0\sqrt{\dfrac{2h}{g}}}
\]

Velocità?
\[
  \vec{V} = \frac{d}{dt}\,\vec{R} = \left(\frac{d}{dt}x, \frac{d}{dt}y\right) = (v_0, -gt)
\]

$t$ finale?
\[
  y(t_{finale}) = y_0 - \frac12 g t^2 = 0 \ \to\ t_{finale} = \sqrt{\frac{2h}{g}}
\]

$v$ finale?
\[
  \vec{V}(t_{finale}) = \left(v_0;\ -\sqrt{2hg}\right)
\]

% ===== da p11.tex =====
% =====================================================================
%  CAPITOLO 10 -- COORDINATE INTRINSECHE   (appunti pp. 32-33)
% =====================================================================
\chapter{Coordinate intrinseche}

Sono delle \evid{coordinate locali} \implica seguono punto per punto la traiettoria.\\
Seguono l'evolversi della traiettoria \implica ascissa curvilinea.

\evid{Definiamo}: \implica \evid{Due versori}:
% disegno: p32-versori (appunti p. 32)
\begin{disegno}
\begin{tikzpicture}[disegno,x=0.12mm,y=-0.12mm]
  \draw[thick,black] (120,275) .. controls (170,195) and (240,172) .. (297,172)
        .. controls (380,170) and (440,232) .. (590,230)
        .. controls (660,228) and (700,210) .. (725,188);
  \draw[vett,accento] (297,172) -- (455,133);
  \node[accento] at (485,128) {\large$\hat{\tau}$};
  \draw[vett,materia!75] (300,174) -- (320,268);
  \node[materia!75] at (362,250) {\large$\vec{n}$};
  \fill[accento] (297,172) circle (2pt);
\end{tikzpicture}
\end{disegno}
\begin{itemize}
  \item \evid{$\hat{\tau}$}: versore \evid{tangente alla traiettoria}
  \item \evid{$\hat{n}$}: versore \evid{normale alla traiettoria}
\end{itemize}

Tali che:
\begin{itemize}
  \item[\implica] tra di loro: \ $\underbrace{\hat{\tau}\cdot\hat{n}}_{\text{prodotto scalare}} = 0 \iff \hat{\tau} \perp \hat{n}$ \quad sono ortogonali
\end{itemize}

\section{Velocità}
\evid{Velocità}:
\begin{formula}
\underbrace{\vec{V}(t)}_{\text{vettore}} = \underbrace{v(t)}_{\text{scalare}}\,\underbrace{\hat{\tau}}_{\text{versore tangente}}
\end{formula}
\begin{significato}
La velocità è sempre tangente alla traiettoria: in queste coordinate basta un numero,
$v(t)$, per descriverla, e la direzione è data da $\hat{\tau}$ punto per punto.
\end{significato}
\begin{itemize}
  \item tang.: \ $(v(t)\hat{\tau})\cdot\hat{\tau} = v(t)$ \qquad ($\hat{\tau}\cdot\hat{\tau} = 1$ perché paralleli)
  \item norm.: \ $(v(t)\hat{\tau})\cdot\hat{n} = 0$ \qquad ($\hat{\tau}\cdot\hat{n} = 0$ perché ortogonali)
\end{itemize}

\section{Accelerazione}
\evid{Accelerazione}
\[
  \vec{a} = \frac{d\vec{V}}{dt} \ \to\ \vec{V} = v\hat{\tau} \ \to\
  \evid{\vec{a} = \underbrace{\dfrac{dv}{dt}}_{\text{diretto come } \hat{\tau}}\hat{\tau} + v\,\underbrace{\dfrac{d\hat{\tau}}{dt}}_{\text{?}}}
\]
Che verso e modulo ha $\dfrac{d\hat{\tau}}{dt}$?
\[
  \hat{\tau} \perp \frac{d\hat{\tau}}{dt} \ \Rightarrow\ \hat{\tau}\,\frac{d\hat{\tau}}{dt} = 0
\]

\begin{dimostrazione}
$\hat{\tau}\cdot\hat{\tau} = |\hat{\tau}|^2 = 1$ \implica facciamo la derivata rispetto al tempo:
\[
  \frac{d\hat{\tau}}{dt}\,\hat{\tau} + \hat{\tau}\,\frac{d\hat{\tau}}{dt}
  = \underbrace{\frac{d|\hat{\tau}|^2}{dt}}_{= 0} \qquad \left(\frac{d}{dt}\,1 = 0\right)
\]
\[
  2\hat{\tau}\,\frac{d\hat{\tau}}{dt} = 0 \iff \evid{\hat{\tau} \perp \dfrac{d\hat{\tau}}{dt}}
\]
\end{dimostrazione}

Ora \chiave{proiettiamo $\vec{a}$ su $\hat{\tau}$} (componente tangenziale):
\[
  \vec{a}\cdot\hat{\tau} = \frac{dv}{dt} + v(t)\underbrace{\left(\hat{\tau}\,\frac{d\hat{\tau}}{dt}\right)}_{= 0}
\]
Dunque \evid{la componente tang.~dell'accelerazione è rappresentata solo da $\dfrac{dv}{dt}$}: come cambia la velocità.

\section{Derivata del versore tangente}
Cioè $\hat{\tau}$ dipende dal tempo ma è legato al punto sulla traiettoria.

Ma\dots\ $\hat{\tau}$ è una funzione di $t$ attraverso $s$:
\evid{$\hat{\tau} = \hat{\tau}(s(t))$} \implica funzione composta \implica deriviamola (derivazione di una composta).
\[
  \frac{d\hat{\tau}}{dt} = \frac{d\hat{\tau}}{ds}\cdot\underbrace{\frac{ds}{dt}}_{= v(t)} = \frac{d\hat{\tau}}{ds}\,v
\]
Ora\dots
\[
  \vec{a} = \frac{dv}{dt}\,\hat{\tau} + v\,\frac{d\hat{\tau}}{dt}
  \qquad
  \evid{v\,\dfrac{d\hat{\tau}}{dt} = v^2\,\dfrac{d\hat{\tau}}{ds}}
\]
\begin{itemize}[nosep]
  \item $v^2$: ci dice come percorriamo la curva (cinematica)
  \item $\dfrac{d\hat{\tau}}{ds}$: ci dice come è fatta la curva (geometria)
\end{itemize}

Che direzione ha? Dato che \evid{$\dfrac{d\hat{\tau}}{ds}\,v = \dfrac{d\hat{\tau}}{dt}$} e $v$ = scalare
\implica hanno la stessa direzione:
\[
  \frac{d\hat{\tau}}{dt} \perp \hat{\tau} \ \Rightarrow\ \frac{d\hat{\tau}}{ds} \perp \hat{\tau}
  \ \Rightarrow\ \evid{\dfrac{d\hat{\tau}}{ds} \parallel \hat{n}}
\]

Analisi dimensionale:
\[
  \left[\frac{d\hat{\tau}}{ds}\right] = \frac{[d\hat{\tau}]}{\underbrace{[ds]}_{\text{lunghezza}}} = \frac{1}{[L]} = [L]^{-1}
  \ \longmapsto\
  \evid{\dfrac{d\hat{\tau}}{ds} = \dfrac{\hat{n}}{\rho}}
  \qquad [\rho] = [L]
\]
Spiegazione: $\dfrac{d\hat{\tau}}{ds}$
\begin{itemize}[nosep]
  \item ha modulo di dim.\ $\left[\frac{1}{L}\right]$, per convenzione lo chiamiamo $\dfrac{1}{\rho}$
  \item ha direzione $\hat{n}$
\end{itemize}
\[
  \frac{d\hat{\tau}}{dt} = \frac{1}{\rho}\,\hat{n}
\]
\evid{$\rho$: raggio di curvatura della traiettoria}.

Se\dots
\begin{itemize}[nosep]
  \item[\implica] curva $=$ circonferenza \implica $\rho$ $=$ raggio
  \item[\implica] curva $=$ retta \implica $\rho$ $=$ infinito
\end{itemize}

\begin{delucidazione}
Tanto più una circonferenza si allarga, tanto più il centro ``$C$'' si allontana, fino a quando
la circonf.\ non collassa in una retta, e il centro dista $\infty$.
\end{delucidazione}

Dunque\dots
\begin{formula}
\vec{a} = \left(\frac{dv}{dt}\right)\hat{\tau} + \frac{v^2}{\rho}\,\hat{n}
\end{formula}
\begin{significato}
L'accelerazione ha due pezzi: $\frac{dv}{dt}$ lungo la traiettoria (si va più o meno veloci)
e $\frac{v^2}{\rho}$ verso il centro della curva (si cambia direzione). Più si va veloci o più
la curva è stretta ($\rho$ piccolo), più grande è l'accelerazione verso l'interno.
\end{significato}
% disegno: p33-curva (appunti p. 33)
\begin{disegno}
\begin{tikzpicture}[disegno,x=0.15mm,y=0.15mm]
  \coordinate (C) at (0,0);
  \def\R{205}\def\a{38}
  \coordinate (P) at (\a:\R);
  \draw[thick,black] (100:\R+12) arc[start angle=100,end angle=-18,radius=\R+12];
  \draw[very thick,nota] (C) -- ($(P)+(\a:12)$);
  \coordinate (P) at (\a:\R+12);
  \draw[green!60!black,thick] ($(P)+(\a+180:14)$) -- ++(\a-90:14) -- ($(P)+(\a-90:14)$);
  \draw[vett,accento] ($(P)+(\a+90:9)$) -- ++(\a+180:110);
  \node[accento] at ($(P)+(\a+180:85)+(\a+90:22)$) {$\hat{n}$};
  \draw[vett,materia!75] (P) -- ++(\a-90:100);
  \node[materia!75] at ($(P)+(\a-90:80)+(\a:24)$) {$\hat{\tau}$};
  \fill (C) circle (1.8pt) node[left] {$C$};
  \node[nota] at ($(C)+(\a:85)+(\a-90:22)$) {\large$\rho$};
\end{tikzpicture}
\end{disegno}
\begin{itemize}
  \item componente tangenziale: $\dfrac{dv}{dt}$
  \item componente normale: sempre $\neq 0$, tranne nel MRU dove $\rho = \infty$
\end{itemize}

% ===== da p12.tex =====
% =====================================================================
%  CAPITOLO 11 -- COORDINATE POLARI   (appunti pp. 34-35)
% =====================================================================
\chapter{Coordinate polari (piane o cilindriche)}

Coordinate polari (piane o cilindriche, dette anche):
in due dimensioni sono polari, in 3D sono cilindriche.

% disegno: p34-cilindro (appunti p. 34)
\begin{disegno}
\begin{tikzpicture}[disegno,x=0.095mm,y=-0.095mm]
  % ---------- cilindro 3D ----------
  \coordinate (O) at (232,295);
  \draw[asse,black] (O) -- (237,40);
  \draw[asse,black] (O) -- (462,287) node[below right] {$y$};
  \draw[asse,black] (O) -- (110,432) node[right=3pt] {$x$};
  \draw[thick,black] (243,152) ellipse[x radius=96,y radius=42];
  \draw[thick,black,dashed] (150,178) -- (163,352);
  \draw[thick,black,dashed] (339,170) -- (340,365);
  \draw[thick,black,dashed] (176,345) .. controls (178,375) and (190,388) .. (240,388)
        .. controls (295,388) and (325,380) .. (333,368);
  \draw[very thick,materia!75] (238,147) -- (283,190);
  \fill (238,147) circle (1.6pt);
  \draw[accento,dashed,thick] (283,190) -- (O);
  \draw[<->,thick,black] (402,125) -- (402,275);
  \node at (440,205) {$z$};
  % ---------- graffa e testo ----------
  \draw[thick,black] (617,105) .. controls (600,115) and (600,160) .. (600,180)
        .. controls (600,195) and (596,200) .. (590,203)
        .. controls (596,206) and (600,215) .. (600,240)
        .. controls (600,300) and (590,340) .. (575,355);
  \node[anchor=west,align=left] at (640,200) {In\\proiezione};
  % ---------- polare 2D ----------
  \coordinate (Q) at (1172,328);
  \draw[asse,black] (1112,328) -- (1468,326) node[below right] {$x$};
  \draw[asse,black] (1171,372) -- (1172,115) node[above right] {$y$};
  \draw[thick,black] (Q) circle[radius=138];
  \coordinate (P) at ($(Q)+(-37:138)$);
  \draw[very thick,materia!75] (Q) -- (P);
  \fill (P) circle (1.8pt);
  \node[above right] at (P) {$P$};
  \node[materia!75] at (1203,262) {\large$\rho$};
  \draw[thick,accento] ($(Q)+(0:70)$) arc[start angle=0,end angle=-37,radius=70];
  \node[accento] at (1265,305) {$\varphi$};
  \draw[->,black!55,thick] (1220,222) -- (1252,172);
  \node[black!55] at (1290,150) {``rho''};
\end{tikzpicture}
\end{disegno}
In proiezione: \ $P(\rho(t), \varphi(t))$ \implica coordinate locali variabili.

\begin{formula}
\begin{cases}
  x = \rho\cos\varphi \\
  y = \rho\sin\varphi \\
  z = z \ \text{(altezza)}
\end{cases}
\quad\longleftrightarrow\quad
\begin{cases}
  \rho = \sqrt{x^2 + y^2} \\
  \varphi = \arctan\!\left(\frac{y}{x}\right) \\
  z = z
\end{cases}
\end{formula}
\begin{significato}
Invece di dire ``quanto a destra e quanto in alto'' ($x, y$), si dice ``quanto lontano dal centro''
($\rho$) e ``in che direzione'' ($\varphi$). Le formule permettono di passare da una descrizione all'altra.
\end{significato}

A cosa servono? \implica Tutti quei \evid{sistemi dove il moto dipende da una determinata origine}
\evid{(es.~moto gravitazionale)}.

\section{Versori polari}
Definiamo i \evid{versori $\hat{\rho}$ e $\hat{\varphi}$} locali.
% disegno: p34-versori (appunti p. 34)
\begin{disegno}
\begin{tikzpicture}[disegno,x=0.1mm,y=-0.1mm]
  \coordinate (O) at (572,308);
  \draw[asse,black] (425,308) -- (868,300) node[below right] {$x$};
  \draw[asse,black] (573,350) -- (573,100) node[left] {$y$};
  % OP e angolo
  \draw[->,thick,black] (O) -- (755,192) node[above right] {$P$};
  \draw[thick,black] (672,306) arc[start angle=0,end angle=-31,radius=100];
  \node at (712,262) {$\varphi$};
  % versori
  \draw[vett,materia!75] (O) -- (645,262);
  \draw[vett,accento] (O) -- (538,252);
  \node[materia!75] at (632,210) {\large$\hat{\rho}$};
  \node[accento] at (518,208) {\large$\hat{\varphi}$};
  \draw[black!55,thick] (632,212) ellipse[x radius=30,y radius=44];
  \draw[black!55,thick] (520,210) ellipse[x radius=38,y radius=40];
  % commenti
  \draw[->,black!55,thick] (488,222) .. controls (440,250) and (405,290) .. (400,342);
  \node at (300,400) {ortogonale a $\hat{\rho}$};
  \draw[->,black!55,thick] (618,256) .. controls (612,300) and (620,345) .. (625,385);
  \node[align=center,anchor=north west] at (605,400) {cambia in\\funzione del\\tempo\\$\Downarrow$\\è una funzione\\di $t$};
  \draw[->,black!55,thick] (710,210) -- (722,152);
  \node[anchor=west] at (690,125) {$\vec{\rho}(t)=\rho(t)\,\hat{\rho}$};
\end{tikzpicture}
\end{disegno}
\[
  \vec{\rho}(t) = \rho(t)\,\hat{\rho}
\]
$\hat{\varphi} \perp \hat{\rho}$ \implica anche nelle rotazioni mantengono l'ortogonalità.\\
$\hat{\varphi}$: ortogonale a $\hat{\rho}$; \ $\hat{\rho}$ cambia in funzione del tempo \implica è una funzione di $t$.

\section{Derivate temporali dei versori polari}
\evid{Derivate temporali dei versori polari}
% disegno: p34-derivate (appunti p. 34)
\begin{disegno}
\begin{tikzpicture}[disegno]
  % assi (senza etichette, come nell'originale)
  \draw[asse] (-1.6,0) -- (1.6,0);
  \draw[asse] (0,-0.8) -- (0,3.0);
  % versori rho(t) e rho(t+dt)
  \draw[vett,thick] (0,0) -- (44:2.1) coordinate (R1);
  \draw[vett,thick] (0,0) -- (58:2.1) coordinate (R2);
  % versori phi(t) e phi(t+dt)
  \draw[vett,thick] (0,0) -- (135:2.1) coordinate (F1);
  \draw[vett,thick] (0,0) -- (148:2.1) coordinate (F2);
  % angolo phi
  \draw (0.55,0) arc (0:44:0.55);
  \node at (22:0.8) {$\varphi$};
  % d rho cappello (blu)
  \draw[->,thick,materia] (R1) -- ($(R1)!0.92!(R2)$);
  \node[materia,anchor=south west] at ($(R1)+(0.0,0.05)$) {$d\hat{\rho}$};
  % d phi a destra (giallo)
  \draw[thick,nota] (40:1.25) arc (40:62:1.25);
  \draw[->,thick,nota] (51:1.25) -- ++(1.75,-0.05) node[right] {$d\varphi$};
  % d phi a sinistra (giallo)
  \draw[thick,nota] (131:1.3) arc (131:152:1.3);
  \node[nota] at (-0.95,1.95) {$d\varphi$};
  % etichette grigie con frecce
  \node[black!55] (Lr) at (1.05,2.8) {$\hat{\rho}(t+dt)$};
  \draw[->,black!55,thick] ($(R2)+(-0.3,-0.25)$) to[bend left=10] ($(R2)+(-0.35,0.75)$);
  \node[black!55] (Lf) at (-1.55,2.35) {$\hat{\varphi}(t+dt)$};
  \draw[->,black!55,thick] ($(F2)+(0.05,0.1)$) to[bend left=15] ($(F2)+(-0.05,0.95)$);
\end{tikzpicture}
\end{disegno}
\[
  \vec{V} = \frac{d\vec{\rho}}{dt} = \frac{d}{dt}\left(\rho(t)\,\hat{\rho}\right)
  = \underbrace{\frac{d\rho(t)}{dt}}_{= \dot{\rho}}\cdot\hat{\rho} + \frac{d\hat{\rho}}{dt}\,\rho(t)
  = \dot{\rho}\,\hat{\rho} + \frac{d\hat{\rho}}{dt}\,\rho(t)
\]
($\rho(t) = \rho$ d'ora in poi)
\[
  |\hat{\rho}(t)| = 1 \qquad |\hat{\varphi}(t)| = 1 \qquad |\hat{\rho}(t + dt)| = 1
\]
\[
  \hat{\rho}(t + dt) = \hat{\rho}(t) + d\hat{\rho} \ \Rightarrow\ d\hat{\rho} = \hat{\rho}(t + dt) - \hat{\rho}(t)
\]
\[
  |d\hat{\rho}| = \underbrace{|\hat{\rho}(t)|}_{= 1}\,d\varphi = d\varphi
\]
Sviluppo il quadrato:
\[
  \underbrace{|\hat{\rho}(t + dt)|^2}_{= 1} = |\hat{\rho}(t) + d\hat{\rho}|^2
  = \underbrace{|\hat{\rho}(t)|^2}_{= 1} + |d\hat{\rho}|^2 + 2\,d\hat{\rho}\cdot\hat{\rho}(t)
\]
\[
  2\,d\hat{\rho}\,\hat{\rho}(t) = 0
  \qquad
  d\hat{\rho} \perp \hat{\rho}(t) \Rightarrow d\hat{\rho} \parallel \hat{\varphi}
\]
\begin{formula}
d\hat{\rho} = d\varphi\,\hat{\varphi}
\qquad
\frac{d\hat{\rho}}{dt} = \hat{\varphi}\,\frac{d\varphi}{dt} = \hat{\varphi}\,\dot{\varphi}
\end{formula}
\begin{significato}
Quando il punto gira di un angolo $d\varphi$, il versore $\hat{\rho}$ ruota con lui: cambia solo
direzione (non lunghezza), e la variazione è perpendicolare a lui, cioè lungo $\hat{\varphi}$.
Più velocemente gira ($\dot{\varphi}$), più rapidamente cambia $\hat{\rho}$.
\end{significato}

\section{Velocità in coordinate polari}
\begin{formula}
\vec{V} = \dot{\rho}\,\hat{\rho} + \rho\,\dot{\varphi}\,\hat{\varphi}
\end{formula}
\begin{itemize}[nosep]
  \item $\dot{\rho}\,\hat{\rho}$: variazione del modulo della traiettoria
  \item $\rho\dot{\varphi}\,\hat{\varphi}$: rotazione della traiettoria
\end{itemize}
\begin{significato}
La velocità ha due parti: quanto ci si allontana o avvicina al centro ($\dot{\rho}$, in direzione radiale)
e quanto si gira attorno al centro ($\rho\dot{\varphi}$, in direzione perpendicolare).
\end{significato}

Interpretazione: \evid{sovrapposizione di un moto rettilineo (direzione $\hat{\rho}$) e una rotazione
rispetto ad esso (direzione $\hat{\varphi}$)}.

\section{Accelerazione in coordinate polari}
\evid{Accelerazione}
\[
  \vec{a} = \frac{d\vec{V}}{dt} = \frac{d}{dt}\left(\dot{\rho}\,\hat{\rho} + \rho\dot{\varphi}\,\hat{\varphi}\right)
  = \ddot{\rho}\,\hat{\rho} + \dot{\rho}\,\underbrace{\frac{d\hat{\rho}}{dt}}_{\dot{\varphi}\hat{\varphi}}
  + \dot{\rho}\dot{\varphi}\,\hat{\varphi} + \rho\ddot{\varphi}\,\hat{\varphi} + \rho\dot{\varphi}\,\frac{d\hat{\varphi}}{dt}
\]
% disegno: p35-dphi (appunti p. 35)
\begin{disegno}
\begin{tikzpicture}[disegno]
  % assi (senza etichette, come nell'originale)
  \draw[asse] (-1.6,0) -- (1.6,0);
  \draw[asse] (0,-0.8) -- (0,3.0);
  % versori rho(t) e rho(t+dt)
  \draw[vett,thick] (0,0) -- (44:2.1) coordinate (R1);
  \draw[vett,thick] (0,0) -- (58:2.1) coordinate (R2);
  % versori phi(t) e phi(t+dt)
  \draw[vett,thick] (0,0) -- (135:2.1) coordinate (F1);
  \draw[vett,thick] (0,0) -- (148:2.1) coordinate (F2);
  % angolo phi
  \draw (0.55,0) arc (0:44:0.55);
  \node at (22:0.8) {$\varphi$};
  % d rho cappello (blu)
  \draw[->,thick,materia] (R1) -- ($(R1)!0.92!(R2)$);
  \node[materia,anchor=south west] at ($(R1)+(0.0,0.05)$) {$d\hat{\rho}$};
  % d phi a destra (giallo)
  \draw[thick,nota] (40:1.25) arc (40:62:1.25);
  \draw[->,thick,nota] (51:1.25) -- ++(1.75,-0.05) node[right] {$d\varphi$};
  % d phi a sinistra (giallo)
  \draw[thick,nota] (131:1.3) arc (131:152:1.3);
  \node[nota] at (-0.95,1.95) {$d\varphi$};
  % etichette grigie con frecce
  \node[black!55] (Lr) at (1.05,2.8) {$\hat{\rho}(t+dt)$};
  \draw[->,black!55,thick] ($(R2)+(-0.3,-0.25)$) to[bend left=10] ($(R2)+(-0.35,0.75)$);
  \node[black!55] (Lf) at (-1.55,2.35) {$\hat{\varphi}(t+dt)$};
  \draw[->,black!55,thick] ($(F2)+(0.05,0.1)$) to[bend left=15] ($(F2)+(-0.05,0.95)$);
  % d phi cappello (blu)
  \draw[->,thick,materia] (F1) -- ($(F1)!0.88!(F2)$);
  \draw[->,thick,materia] ($(F1)+(-0.05,0.05)$) to[out=170,in=90,looseness=0.9] (-2.6,-0.15);
  \node[materia,below] at (-2.5,-0.2) {$d\hat{\varphi}$};
\end{tikzpicture}
\end{disegno}
\[
  \hat{\varphi}(t + dt) - \hat{\varphi}(t) = d\hat{\varphi}
  \ \longmapsto\
  |d\hat{\varphi}| = d\varphi\,\underbrace{|\hat{\varphi}(t)|}_{= 1} = d\varphi
\]
\[
  \frac{d\hat{\varphi}}{dt} = \frac{d\varphi}{dt}\,(-\hat{\rho}) = -\hat{\rho}\,\dot{\varphi}
\]
\[
  \vec{a} = \ddot{\rho}\,\hat{\rho} + \underbrace{\dot{\rho}\dot{\varphi}\,\hat{\varphi} + \dot{\rho}\dot{\varphi}\,\hat{\varphi}}_{2\dot{\rho}\dot{\varphi}\hat{\varphi}}
  + \rho\ddot{\varphi}\,\hat{\varphi} + \rho\dot{\varphi}\,\frac{d\hat{\varphi}}{dt}
  \qquad
  \rho\dot{\varphi}\left(-\hat{\rho}\dot{\varphi}\right) = -\rho\dot{\varphi}^2\,\hat{\rho}
\]
Dunque\dots
\begin{formula}
\vec{a} = \hat{\rho}\left(\ddot{\rho} - \rho\dot{\varphi}^2\right) + \hat{\varphi}\left(\rho\ddot{\varphi} + 2\dot{\rho}\dot{\varphi}\right)
\end{formula}
\begin{itemize}[nosep]
  \item $\hat{\rho}(\dots)$: radiale; \ $\ddot{\rho}$ in direzione di $\hat{\rho}$; \ $-\rho\dot{\varphi}^2$: centripeta ($-\hat{\rho}$)
  \item $\hat{\varphi}(\dots)$: tangenziale; \ $\rho\ddot{\varphi}$: accelerazione tangenziale;
        \ $2\dot{\rho}\dot{\varphi}$: termine misto, moto vario
\end{itemize}
\begin{significato}
Lungo il raggio: $\ddot{\rho}$ è l'accelerazione con cui ci si allontana dal centro, mentre
$-\rho\dot{\varphi}^2$ è la parte centripeta dovuta al girare. In direzione perpendicolare:
$\rho\ddot{\varphi}$ dice se si gira sempre più in fretta, e $2\dot{\rho}\dot{\varphi}$ compare
quando ci si allontana dal centro mentre si gira (si ``sente'' quando si cammina verso il bordo di una giostra).
\end{significato}

\section{Derivate dei versori polari usando la rappres.\ cartesiana}
\evid{Derivate dei versori polari usando la rappres.~cartesiana}
% disegno: p35-cartesiana (appunti p. 35)
\begin{disegno}
\begin{tikzpicture}[disegno,scale=1.1]
  \def\a{25}
  % assi
  \draw[asse] (-1.4,0) -- (3.2,0) node[below right=-1pt] {$x$};
  \draw[asse] (0,-0.6) -- (0,2.1) node[right=1pt] {$y$};
  % versori
  \coordinate (R) at (\a:2.0);
  \coordinate (F) at (90+\a:2.0);
  \draw[vett,thick] (0,0) -- (R);
  \draw[vett,thick] (0,0) -- (F);
  % proiezioni
  \draw[tratt] (R) -- (R|-0,0);
  \draw[tratt] (R) -- (R-|0,0);
  \draw[tratt] (F) -- (F|-0,0);
  \draw[tratt] (F) -- (F-|0,0);
  \draw[->,thick] ($(R|-0,0)+(-0.3,0)$) -- (R|-0,0);
  % angoli
  \draw (0.75,0) arc (0:\a:0.75);
  \node at (\a/2:1.0) {$\varphi$};
  \draw (0,0.55) arc (90:90+\a:0.55);
  \node at (90+\a/2:0.95) {$\varphi$};
  % etichette
  \draw ($(R)+(0.05,0.05)$) -- ($(R)+(0.3,0.45)$);
  \node at ($(R)+(0.45,0.7)$) {$\hat{\rho}$};
  \node at ($(F)+(-0.2,0.4)$) {$\hat{\varphi}$};
\end{tikzpicture}
\end{disegno}
\[
  \hat{\rho} = (\cos\varphi, \sin\varphi)
  \qquad
  \hat{\varphi} = (-\sin\varphi, \cos\varphi)
  \qquad
  \hat{\rho}\cdot\hat{\varphi} = \cos\varphi(-\sin\varphi) + \sin\varphi\cdot\cos\varphi = 0
\]
\[
  \begin{cases}
    \dfrac{d\hat{\rho}}{dt} = \dfrac{d}{dt}(\cos\varphi, \sin\varphi)
      = \Big(\underbrace{\tfrac{d}{d\varphi}(\cos\varphi)}_{-\sin\varphi}\underbrace{\tfrac{d\varphi}{dt}}_{\dot{\varphi}},\
             \underbrace{\tfrac{d}{d\varphi}(\sin\varphi)}_{\cos\varphi}\underbrace{\tfrac{d\varphi}{dt}}_{\dot{\varphi}}\Big)
      = \dot{\varphi}(-\sin\varphi, \cos\varphi) = \dot{\varphi}\,\hat{\varphi} \\[4ex]
    \dfrac{d\hat{\varphi}}{dt} = \dfrac{d}{dt}(-\sin\varphi, \cos\varphi)
      = \Big(\underbrace{\tfrac{d}{d\varphi}(-\sin\varphi)}_{-\cos\varphi}\underbrace{\tfrac{d\varphi}{dt}}_{\dot{\varphi}},\
             \underbrace{\tfrac{d(\cos\varphi)}{d\varphi}}_{-\sin\varphi}\underbrace{\tfrac{d\varphi}{dt}}_{\dot{\varphi}}\Big)
      = -\dot{\varphi}(\cos\varphi, \sin\varphi) = -\dot{\varphi}\,\hat{\rho}
  \end{cases}
\]

% ===== da p13.tex =====
% ---------------------------------------------------------------------
%  (segue CAPITOLO 11)  Velocità e accelerazione angolare (appunti p. 36)
% ---------------------------------------------------------------------
\section{Velocità angolare \texorpdfstring{$\vec{\omega}$}{omega}}
\[
  |\vec{\omega}| = |\dot{\varphi}| \qquad \omega = \dot{\varphi} \ \text{(con il suo segno)}
\]
\begin{formula}
\underbrace{\vec{\omega} = \omega\,\hat{z}}_{\text{ortogonale al piano del moto}} \qquad \text{(moto nel piano } (x, y))
\end{formula}
\begin{significato}
La velocità angolare dice quanto velocemente cambia l'angolo: è un vettore perpendicolare
al piano in cui si gira, e il suo verso indica il senso di rotazione (orario o antiorario).
\end{significato}

\section{Versori degli assi polari}
\evid{Versori degli assi polari} (prodotto vettoriale):
% disegno: p36-assi (appunti p. 36)
\begin{disegno}
\begin{tikzpicture}[disegno,scale=1.3]
  % assi
  \draw[asse] (0,0) -- (0,1.9) node[above right=-1pt] {$z$};
  \draw[asse] (0,0) -- (2.2,0) node[right] {$x$};
  \draw[asse] (0,0) -- (-1.2,-1.0) node[below right=-1pt] {$y$};
  % piano xz (blu): tratti verticali
  \foreach \u in {0.6,1.15,1.6,2.0}
    \draw[thick,materia] (\u,0.05) -- (\u,1.45-0.05*\u);
  % piano yz (rosso): tratti paralleli a y
  \foreach \v in {0.25,0.65,1.05,1.45}
    \draw[thick,accento] (-0.25,\v) -- ++(-0.85,-0.65);
  % piano xy (verde): tratti orizzontali
  \foreach \s/\l in {0.3/1.0,0.55/1.4,0.8/1.8,1.05/1.75}
    \draw[thick,esempio] (-0.85*\s+0.55,-0.75*\s-0.05) -- ++(\l,-0.25);
\end{tikzpicture}
\end{disegno}
\begin{formula}
\begin{cases}
  \hat{\rho}\wedge\hat{\varphi} = \hat{z} \\
  \hat{z}\wedge\hat{\rho} = \hat{\varphi} \\
  \hat{\varphi}\wedge\hat{z} = \hat{\rho}
\end{cases}
\end{formula}

\evid{Velocità}
\[
  \vec{V} = \dot{\rho}\,\hat{\rho} + \rho\dot{\varphi}\,\hat{\varphi}
  \ \longmapsto\
  \rho\dot{\varphi}\,\hat{\varphi} = \rho\omega\,(\hat{z}\wedge\hat{\rho})
  = \underbrace{(\omega\hat{z})}_{= \vec{\omega}}\wedge\underbrace{(\rho\hat{\rho})}_{= \vec{\rho}}
\]
\begin{formula}
\vec{V} = \dot{\rho}\,\hat{\rho} + \vec{\omega}\wedge\vec{\rho}
\end{formula}
\begin{significato}
La parte di velocità dovuta alla rotazione si scrive come $\vec{\omega}\wedge\vec{\rho}$:
è perpendicolare al raggio e vale $\omega\rho$, cioè più si è lontani dal centro più si va veloci
a parità di velocità angolare.
\end{significato}

\section{Accelerazione angolare}
\evid{Accelerazione angolare}
\begin{formula}
\vec{\alpha} = \frac{d\vec{\omega}}{dt} = \dot{\omega}\,\hat{z} = \ddot{\varphi}\,\hat{z}
\end{formula}
\begin{significato}
Dice quanto rapidamente cambia la velocità angolare: se è nulla si gira a ritmo costante.
\end{significato}
$\vec{\alpha} \parallel \vec{\omega}$: non sono necessariamente equiversi.

\[
  \vec{a} = \hat{\rho}\left(\ddot{\rho} - \rho\dot{\varphi}^2\right) + \hat{\varphi}\left(2\dot{\rho}\dot{\varphi} + \rho\ddot{\varphi}\right)
  \qquad \text{Utilizziamo } \vec{\omega} \text{ e } \vec{\alpha}
\]
\begin{itemize}[nosep]
  \item $\hat{\rho}\ddot{\rho}$: invariato
  \item $-\hat{\rho}\rho\underbrace{\dot{\varphi}^2}_{\omega^2} = -\omega^2\underbrace{(\rho\hat{\rho})}_{\vec{\rho}} = -\omega^2\vec{\rho}$
  \item $\hat{\varphi}\rho\ddot{\varphi} = (\hat{z}\wedge\hat{\rho})\rho\ddot{\varphi}
         = \underbrace{(\hat{z}\ddot{\varphi})}_{\vec{\alpha}}\wedge\underbrace{(\hat{\rho}\rho)}_{\vec{\rho}} = \vec{\alpha}\wedge\vec{\rho}$
  \item $2\dot{\rho}\dot{\varphi}\hat{\varphi} = 2\dot{\rho}\underbrace{\omega(\hat{z}}_{\vec{\omega}}\wedge\hat{\rho}) = 2\dot{\rho}(\vec{\omega}\wedge\hat{\rho})$
\end{itemize}
\begin{formula}
\vec{a} = \hat{\rho}\,\ddot{\rho} - \omega^2\vec{\rho} + \vec{\alpha}\wedge\vec{\rho} + 2\dot{\rho}\,(\vec{\omega}\wedge\vec{\rho})
\end{formula}
\begin{significato}
È la stessa accelerazione in coordinate polari, scritta con $\vec{\omega}$ e $\vec{\alpha}$:
$-\omega^2\vec{\rho}$ è la parte centripeta, $\vec{\alpha}\wedge\vec{\rho}$ la parte tangenziale dovuta
al cambiare della rotazione, l'ultimo termine compare solo se la distanza dal centro cambia mentre si gira.
\end{significato}

% =====================================================================
%  CAPITOLO 12 -- MOTO CIRCOLARE   (appunti pp. 37-...)
% =====================================================================
\chapter{Moto circolare}

\evid{La traiettoria è una circ.~di raggio $R$}.\\
C'è \evid{una sola legge oraria} \implica $\varphi(t)$.
% disegno: p37-circ (appunti p. 37)
\begin{disegno}
\begin{tikzpicture}[disegno]
  \draw[asse] (-1.6,0) -- (1.9,0);
  \draw[asse] (0,-1.0) -- (0,1.7);
  \draw[thick] (0,0) ellipse (1.1 and 0.85);
  \coordinate (P) at (40:1.0);
  \draw[thick] (0,0) -- (P);
  \fill (P) circle (1.5pt);
  \node[above right=-1pt] at (P) {$P$};
  \node at (0.25,0.5) {$R$};
  \draw (0.5,0) arc (0:36:0.5);
  \draw (22:0.55) -- (1.45,0.45) node[right] {$\varphi(t)$};
\end{tikzpicture}
\end{disegno}

\begin{delucidazione}
È un moto \evid{unidimensionale} con \evid{una sola legge oraria} $\varphi(t)$.
\end{delucidazione}

\section{Velocità e accelerazione}
\evid{Velocità e accelerazione}

\evid{Formule generali}, \ $\dot{\rho} = \ddot{\rho} = 0$ \implica
\begin{formula}
\begin{cases}
  \vec{V} = \rho\dot{\varphi}\,\hat{\varphi} = \vec{\omega}\wedge\vec{\rho} \\
  \vec{a} = \underbrace{-\omega^2\vec{\rho}}_{\text{centripeta}} + \underbrace{\vec{\alpha}\wedge\vec{\rho}}_{\text{tangente}}
          = -\rho\dot{\varphi}^2\,\hat{\rho} + \rho\ddot{\varphi}\,\hat{\varphi}
\end{cases}
\end{formula}
\begin{significato}
Nel moto circolare la velocità è sempre tangente alla circonferenza. L'accelerazione ha sempre
una parte verso il centro ($\omega^2\rho$), necessaria per curvare, e una parte tangente solo se
la rotazione accelera o rallenta.
\end{significato}
(La parte tangente nell'uniforme non c'è.)

$\vec{\omega}\wedge\vec{\rho} \perp \vec{\rho}$ \implica velocità $\perp$ raggio \implica $\vec{V} \parallel$ tangente.

% disegno: p37-acc (appunti p. 37)
\begin{disegno}
\begin{tikzpicture}[disegno]
  % commento a sinistra
  \node[align=left,anchor=east] (S) at (-2.35,0.3) {Si suppone\\ che il\\ moto sia\\ \colorbox{evid}{antiorario}};
  \draw[thick] (-2.1,1.0) .. controls (-2.25,1.0) and (-2.15,0.4) .. (-2.3,0.4)
               .. controls (-2.15,0.4) and (-2.25,-0.2) .. (-2.1,-0.2);
  % circonferenza
  \draw[thick] (0,0) circle (1.25);
  \fill (0,0) circle (1.6pt);
  \coordinate (P) at (52:1.25);
  % raggio R
  \draw[thick] (0,0) -- (52:0.75);
  \node at ($(52:0.62)+(-0.12,0.22)$) {$R$};
  % accelerazione centripeta (rosso, verso il centro)
  \draw[->,thick,accento] (P) -- (52:0.72);
  \fill (P) circle (1.8pt);
  % velocita (blu, tangente)
  \draw[->,very thick,materia] (P) -- ($(P)+(142:0.85)$);
  \node[materia] at ($(P)+(0.05,0.85)$) {$\vec{v}$};
  % a tangenziale (rosso, doppia freccia tangente)
  \draw[<->,thick,accento] ($(P)+(0.3,0.75)$) -- ($(P)+(0.42,-0.2)$);
  \node[accento,anchor=west] at ($(P)+(0.45,0.75)$) {$\vec{a}_{\mathrm{tang}}$};
  % omega (verso antiorario)
  \node at (-1.0,1.55) {$\omega$};
  \draw[->,thick] (-0.55,1.45) to[bend right=25] (-1.35,0.95);
  % freccia grigia verso a centrifuga
  \draw[->,thick,black!55] ($(52:0.72)+(0.25,-0.1)$) to[out=-20,in=95] (1.45,-1.5);
  \node[accento,anchor=north west] at (1.25,-1.55) {$\vec{a}_{\text{centrifuga}}$};
  % commento a destra
  \node[align=center,anchor=west] at (2.0,0.45) {Accelerazione\\[1pt] $-\omega^{2}\vec{\rho}$\\[1pt] centrifuga};
\end{tikzpicture}
\end{disegno}
Si suppone che il moto sia antiorario.
\begin{itemize}[nosep]
  \item Accelerazione $-\omega^2\vec{\rho}$: centripeta
  \item $\vec{\alpha}\wedge\vec{\rho}$ ($\parallel \hat{\varphi}$): tangenziale \implica $a_{tangenziale} = 0$ se $\ddot{\varphi} = 0$
  \item $-\omega^2\vec{\rho}$ sempre $\neq 0$
\end{itemize}
\implica \evid{Nel moto circolare esiste sempre un'accelerazione centripeta.}

\subsection{\texorpdfstring{$\vec{a}$}{a} nel moto circolare con il calcolo diretto}
\[
  \vec{a} = \frac{d}{dt}\,\vec{V} = \frac{d}{dt}(\vec{\omega}\wedge\vec{\rho})
  = \underbrace{\frac{d\vec{\omega}}{dt}}_{\vec{\alpha}}\wedge\vec{\rho} + \vec{\omega}\wedge\frac{d\vec{\rho}}{dt}
  = \vec{\alpha}\wedge\vec{\rho} + \vec{\omega}\wedge\vec{V}
  = \vec{\alpha}\wedge\vec{\rho} + \vec{\omega}\wedge(\vec{\omega}\wedge\vec{\rho})
\]
Triplo prodotto vettoriale:
\[
  \evid{\vec{A}\wedge(\vec{B}\wedge\vec{C}) = (\vec{A}\cdot\vec{C})\vec{B} - (\vec{A}\cdot\vec{B})\vec{C}}
\]
\[
  \vec{\omega}\wedge(\vec{\omega}\wedge\vec{\rho}) = \underbrace{(\vec{\omega}\cdot\vec{\rho})}_{= 0}\vec{\omega} - (\omega^2)\vec{\rho} = -\omega^2\vec{\rho}
\]
\begin{formula}
\vec{a} = \vec{\alpha}\wedge\vec{\rho} - \omega^2\vec{\rho}
\end{formula}

\section{Moto circolare uniforme}
\evid{Moto circolare uniforme}: \ $\ddot{\varphi} = \dot{\omega} = 0$, \ \evid{$\vec{\omega}$ costante}.

Formule \implica
\begin{formula}
\vec{V} = \vec{\omega}\wedge\vec{\rho} \qquad \text{e} \qquad \vec{a} = -\omega^2\vec{\rho}
\end{formula}
\begin{significato}
Anche se il modulo della velocità non cambia, la sua direzione cambia continuamente: per questo
c'è sempre un'accelerazione, diretta verso il centro e di modulo $\omega^2\rho$.
\end{significato}

Moto circolare uniforme \evid{è sempre accelerato}.
\[
  \begin{cases}
    \textbf{Periodo}: \ T = \dfrac{2\pi}{\omega} \quad [T] & \text{(tempo per percorrere una circ.)} \\[2ex]
    \textbf{Frequenza}: \ \nu = \dfrac{1}{T} = \dfrac{\omega}{2\pi} \quad \left[\tfrac{1}{T}\right] & \text{(numero di giri per unità di tempo)}
  \end{cases}
\]

% ===== da p14.tex =====
% ---------------------------------------------------------------------
%  (segue CAPITOLO 12)  Moto circolare: coordinate cartesiane e intrinseche (appunti pp. 38-39)
% ---------------------------------------------------------------------
\section{In coordinate cartesiane}
\evid{In coordinate cartesiane}
\[
  \underbrace{\begin{cases}
    x = R\cos[\varphi(t)] \\
    y = R\sin[\varphi(t)]
  \end{cases}}_{\text{posizione (composte)}}
  \Rightarrow
  \underbrace{\begin{cases}
    \dot{x} = -R\sin(\varphi)\,\dot{\varphi}(t) \\
    \dot{y} = R\cos(\varphi)\,\dot{\varphi}(t)
  \end{cases}}_{\text{velocità}}
\]
Accelerazione:
\[
  \Rightarrow
  \begin{cases}
    \ddot{x} = -R\left(\cos(\varphi)\dot{\varphi}^2 + \sin(\varphi)\ddot{\varphi}\right) \\
    \ddot{y} = R\left(-\sin(\varphi)\dot{\varphi}^2 + \cos(\varphi)\ddot{\varphi}\right)
  \end{cases}
  =
  \begin{cases}
    \ddot{x} = -R\sin(\varphi)\ddot{\varphi} \\
    \ddot{y} = R\cos(\varphi)\ddot{\varphi}
  \end{cases}
\]
\begin{delucidazione}
!!! \ Da rivedere (riferito al passaggio qui sopra).
\end{delucidazione}

\section{Nel moto circolare uniforme}
\evid{Nel moto circolare uniforme}
\[
  \ddot{\varphi} = \dot{\omega} = 0 \ \Rightarrow\ \omega = \text{costante} \qquad \varphi = \omega t + \varphi_0
\]
\[
  \underbrace{\begin{cases}
    x = R\cos(\omega t + \varphi_0) \\
    y = R\sin(\omega t + \varphi_0)
  \end{cases}}_{\text{pos.}}
  \Rightarrow
  \underbrace{\begin{cases}
    \dot{x} = -\omega R\sin(\omega t + \varphi_0) \\
    \dot{y} = \omega\cos(\omega t + \varphi_0)
  \end{cases}}_{\text{vel.}}
  \Rightarrow
  \underbrace{\begin{cases}
    \ddot{x} = -\omega^2 x \\
    \ddot{y} = -\omega^2 y
  \end{cases}}_{\text{acc.}}
  \quad \evid{Moti armonici}
\]
\begin{significato}
Guardando solo l'ombra del moto circolare uniforme su un asse ($x$ o $y$) si vede un moto
armonico: avanti e indietro con accelerazione proporzionale allo spostamento e opposta.
\end{significato}

\begin{delucidazione}
Vedere didascalia su slides.
\end{delucidazione}

\begin{nota}
Attenzione! Pulsazione armonica e velocità angolare sono concettualmente distinti:
\begin{itemize}[nosep]
  \item la velocità angolare è una funzione del tempo \implica $\dfrac{d\varphi}{dt}$
  \item la pulsazione non è una funzione del tempo \implica $\dfrac{2\pi}{T}$
\end{itemize}
Nel caso del moto circ.\ uniform.\ si equivalgono.
\end{nota}

\section{Coordinate intrinseche}
\evid{Coordinate intrinseche}
% disegno: p38-intr (appunti p. 38)
\begin{disegno}
\begin{tikzpicture}[disegno,scale=1.2]
  \draw[thick] (0,0) circle (1);
  \coordinate (P) at (0,1);
  \draw[vett] (P) -- (1.6,1) node[right] {$\hat{\tau}$};
  \draw[->,thick] (P) -- (0,0.45);
  \node at (0.28,0.55) {$\vec{n}$};
  \fill (P) circle (1.8pt);
\end{tikzpicture}
\end{disegno}
\[
  \begin{cases}
    \hat{\tau} = \hat{\varphi} \\
    \hat{n} = -\hat{\rho}
  \end{cases}
  \qquad
  \text{velocità:}\quad \vec{V} = \vec{\omega}\wedge\vec{\rho} = \omega|\vec{\rho}|\,\hat{\varphi} = \omega R\,\hat{\tau}
\]
($\vec{\omega}$ e $\vec{\rho}$ ortogonali)

Accelerazione:
\[
  \vec{a} = -\omega^2\vec{\rho} + \vec{\alpha}\wedge\vec{\rho}
  = -R\dot{\varphi}^2\,\hat{\rho} + R\ddot{\varphi}\,\hat{\varphi}
  = \underbrace{-R\left(\frac{v^2}{R^2}\right)\hat{\rho}}_{\hat{n}} + \underbrace{R\dot{\omega}}_{\frac{dv}{dt}}\,\underbrace{\hat{\varphi}}_{\hat{\tau}}
\]
\[
  \ddot{\varphi} = \dot{\omega}
  \qquad
  v = \omega R = R\dot{\varphi} \Rightarrow \dot{\varphi} = \frac{v}{R}
  \qquad
  \frac{dv}{dt} = \frac{d}{dt}(\omega R) = R\,\frac{d\omega}{dt} = R\dot{\omega}
\]
\begin{formula}
\vec{a} = \frac{v^2}{R}\,\hat{n} + \left(\frac{dv}{dt}\right)\hat{\tau}
\end{formula}
\begin{significato}
Nel moto circolare il raggio di curvatura è proprio il raggio $R$: la parte verso il centro vale
$v^2/R$ e c'è sempre; quella tangente vale $dv/dt$ e c'è solo se il modulo della velocità cambia.
\end{significato}

\begin{delucidazione}
!!! \ Vedere slides.
\end{delucidazione}

\section{Esempio: calcolo raggio di curvatura e uso coordinate intrinseche}
\begin{esempio}{grave lanciato orizzontalmente}
% disegno: p39-grave (appunti p. 39)
\begin{disegno}
\begin{tikzpicture}[disegno,materia,scale=0.75]
  \draw[asse,materia] (-0.15,0) -- (3.6,0) node[right] {$x$};
  \draw[asse,materia] (0,-0.25) -- (0,2.9) node[right=1pt] {$y$};
  \draw[curva] (0,2.1) .. controls (1.0,2.1) and (1.9,1.4) .. (2.05,0);
  \fill[materia] (2.05,0) circle (1.8pt);
  \node[below=2pt] at (2.05,0) {$\bar{x}$};
\end{tikzpicture}
\end{disegno}
Raggio di curvatura?
\begin{enumerate}[nosep]
  \item Determiniamo $\vec{V}(t) \to \hat{\tau} = \dfrac{\vec{V}(t)}{|\vec{V}(t)|}$
  \item $\hat{n}\cdot\hat{\tau} = 0$
  \item $\vec{a}(t)$ \ $(-g\hat{y}) \to a_n = \vec{a}\cdot\hat{n}$
  \item $\rho = \dfrac{v^2}{a_n}$
\end{enumerate}

\textbf{(1)}
\[
  \vec{V}(t) = (v_0, -gt) \qquad \hat{\tau} = \frac{(v_0, -gt)}{\sqrt{v_0^2 + g^2t^2}}
\]
\textbf{(2)}
% disegno: p39-versori (appunti p. 39)
\begin{disegno}
\begin{tikzpicture}[disegno,materia,scale=0.8]
  \draw[asse,materia] (-0.3,0) -- (2.9,0) node[below right] {$\hat{x}$};
  \draw[asse,materia] (0,-0.3) -- (0,1.75) node[above right] {$\hat{y}$};
  % traiettoria
  \draw[curva] (0,1.15) .. controls (0.7,1.1) and (1.35,0.6) .. (1.5,-0.15);
  \coordinate (P) at (0.85,0.93);
  \fill (P) circle (1.8pt);
  % versori
  \draw[->,thick] (P) -- ++(-28:1.0) node[right] {$\hat{\tau}$};
  \draw[->,thick] (P) -- ++(-122:0.5);
  % etichetta n con freccia verso il commento
  \node (n) at (0.85,0.18) {$\hat{n}$};
  \draw[thick] (0.68,0.12) .. controls (0.68,0.02) and (1.02,0.02) .. (1.02,0.12);
  \draw[->,thick] (0.85,0.03) -- (0.85,-0.55);
  \node[anchor=north west,align=left] at (0.55,-0.6) {$\hat{n}$ ha entrambe\\ \hspace*{1.2em}comp.\ neg};
\end{tikzpicture}
\end{disegno}
$\hat{n}$ ha entrambe comp.\ neg.:
\[
  \hat{n} = \frac{(-gt, -v_0)}{\sqrt{g^2t^2 + v_0^2}} \ \to\ \hat{n}\cdot\hat{\tau} = 0
\]
(scambio le comp.\ di $\hat{\tau}$ ma metto tutte e due negative)

\textbf{(3)}
\[
  \vec{a} = (0, -g)
  \qquad
  \vec{a}\cdot\hat{n} = a_n = -g\hat{y}\cdot\frac{(\overbrace{-gt\,\hat{x}}^{= 0}, -v_0\,\hat{y})}{\sqrt{g^2t^2 + v_0^2}}
  = \frac{v_0 g}{\sqrt{v_0^2 + g^2t^2}}
\]
\textbf{(4)}
\[
  \rho = \frac{v^2}{a_n} = \frac{(v_0^2 + g^2t^2)}{g v_0}\sqrt{v_0^2 + g^2t^2} = \frac{(v_0^2 + g^2t^2)^{3/2}}{g v_0}
\]
\end{esempio}

% ===== da r2.tex =====
% =====================================================================
%  RIEPILOGO 2 -- MOTI NEL PIANO E NELLO SPAZIO   (pagina aggiunta, non negli appunti)
% =====================================================================
\riepilogo{Moti nel piano e nello spazio}

In più dimensioni posizione, velocità e accelerazione diventano vettori. La velocità è sempre
tangente alla traiettoria; l'accelerazione si divide in una parte tangente (cambia il modulo)
e una normale (cambia la direzione). Le coordinate si scelgono in base al moto.

\begin{mappa}
  \node[nodoevid] (r) {$\vec{R}(t)$};
  \node[nodo, right=1.2cm of r] (v) {$\vec{V} = \dot{\vec{R}}$\\ \small tangente};
  \node[nodo, right=1.2cm of v] (a) {$\vec{a} = \dot{\vec{V}}$};
  \node[nodo, above right=0.1cm and 1.1cm of a] (at) {$a_t = \frac{dv}{dt}$\\ \small modulo};
  \node[nodo, below right=0.1cm and 1.1cm of a] (an) {$a_n = \frac{v^2}{\rho}$\\ \small direzione};
  \draw[freccia] (r) -- (v);
  \draw[freccia] (v) -- (a);
  \draw[freccia] (a.east) -- (at.west);
  \draw[freccia] (a.east) -- (an.west);
  \node[nodoevid, below=1.1cm of v] (coord) {Coordinate};
  \node[nodo, below left=0.5cm and 0.4cm of coord] (cart) {Cartesiane\\ \small proiettili};
  \node[nodo, below=0.5cm of coord] (intr) {Intrinseche\\ \small $\hat{\tau}, \hat{n}$};
  \node[nodo, below right=0.5cm and 0.4cm of coord] (pol) {Polari\\ \small $\hat{\rho}, \hat{\varphi}$\\ \small moto circolare};
  \foreach \n in {cart, intr, pol} \draw[freccia] (coord) -- (\n);
\end{mappa}

\begin{formulario}
\vec{V} = \frac{d\vec{R}}{dt} & Velocità vettoriale, sempre tangente alla traiettoria. \\
\vec{a} = \frac{d\vec{V}}{dt} & Accelerazione vettoriale: cambia se cambia il modulo o la direzione della velocità. \\
\begin{array}{l} x = v_0\cos\theta\, t \\ y = v_0\sin\theta\, t - \frac12 g t^2 \end{array}
  & Proiettile: uniforme in orizzontale, uniformemente accelerato in verticale. \\
h = \frac{v_0^2\sin^2\theta}{2g} & Quota massima: dipende solo dalla velocità verticale iniziale. \\
G = \frac{v_0^2\sin(2\theta)}{g} & Gittata: massima a $45^\circ$, uguale per angoli complementari. \\
y = \frac{v_0^2}{2g} - \frac{gx^2}{2v_0^2} & Parabola di sicurezza: sopra non si viene mai colpiti. \\
\vec{a} = \frac{dv}{dt}\hat{\tau} + \frac{v^2}{\rho}\hat{n} & Coordinate intrinseche: parte tangente (modulo) $+$ parte normale verso il centro di curvatura. \\
\vec{V} = \dot{\rho}\hat{\rho} + \rho\dot{\varphi}\hat{\varphi} & Polari: allontanarsi dal centro $+$ girare attorno al centro. \\
\vec{a} = \hat{\rho}(\ddot{\rho} - \rho\dot{\varphi}^2) + \hat{\varphi}(\rho\ddot{\varphi} + 2\dot{\rho}\dot{\varphi}) & Polari: radiale (con la parte centripeta) $+$ trasversale. \\
\vec{V} = \vec{\omega}\wedge\vec{\rho} & Moto circolare: velocità perpendicolare al raggio, modulo $\omega R$. \\
\vec{a} = \vec{\alpha}\wedge\vec{\rho} - \omega^2\vec{\rho} & Moto circolare: tangenziale (se $\omega$ cambia) $+$ centripeta (sempre). \\
T = \frac{2\pi}{\omega},\ \ \nu = \frac{\omega}{2\pi} & Moto circolare uniforme: tempo di un giro e giri al secondo. \\
\end{formulario}


\end{document}
